\chapter{Plurisubharmonic functions and singularities}
\label{chap:framework}

A function whose complex Hessian is bounded below by a smooth form becomes plurisubharmonic after a local smooth correction. This brings the local theory of plurisubharmonic functions to a compact complex manifold, provided the smooth remainder left by the correction is kept when potentials are glued, when their Hessians are multiplied and when one integrates by parts. Normalization by the supremum then controls the resulting families globally, giving compactness, uniform exponential integrability and capacity estimates on the whole manifold. Products of Hessians require more care than the potentials themselves: without a common bound, convergence of potentials in $L^1$, or even in capacity, does not prevent their Monge--Amp\`ere measures from concentrating.

To keep track of singularities under these operations, we compare potentials in two ways: by domination up to an additive constant, which leads to ordinary singularity types and their envelopes, and by the local integrability of holomorphic functions against exponential weights, which leads to multiplier ideals. Lelong numbers, analytic approximation and modifications are the main tools for computing the latter.

\section{Plurisubharmonic potentials}

Throughout the chapter, $X$ is a connected compact complex manifold of dimension $n\geq1$, $\theta$ is a smooth real $(1,1)$-form and $\omega$ is a Hermitian form on $X$, unless a statement specifies another setting. There are two distinct extensions of the K\"ahler setting: one can study a closed background on a non-K\"ahler manifold, or allow the background itself to be nonclosed. In both cases positivity is expressed through quasi-plurisubharmonic functions.

On an open set $U$ in a complex manifold, a \emph{plurisubharmonic function} is a locally integrable upper semicontinuous function satisfying the submean inequality on every complex line in local coordinates. We write $\PSH(U)$ for the set of these functions.\index{Plurisubharmonic function!local definition} A function is \emph{quasi-plurisubharmonic} if it is locally the sum of a plurisubharmonic function and a smooth real function; the set of such functions is denoted by $\QPSH(U)$. The following relative positivity condition makes sense on any such open set as well.

\begin{definition}\label{def:theta-psh}\index{PSH|see{Plurisubharmonic function}}\index{Quasi-plurisubharmonic function|see{Plurisubharmonic function}}\index{Plurisubharmonic function!relative to a background}
A quasi-plurisubharmonic function $u$ on $X$ is \emph{$\theta$-plurisubharmonic} if
\[
 \theta+\ddc u\geq0
\]
as a current. We denote the set of these functions by $\PSH(X,\theta)$.
\end{definition}

Since $\omega>0$, every constant function belongs to $\PSH(X,\omega)$.

\Needspace{11\baselineskip}
\begin{proposition}\label{prop:local-correction}\index{Local correction}\index{Background form!local correction}
Every real distribution $u$ on $X$ with $\theta+\ddc u\geq0$ has a unique $\theta$-plurisubharmonic representative. On each relatively compact coordinate ball with coordinates $z=(z_1,\ldots,z_n)$, there is a constant $A>0$ such that $u+A|z|^2$ is plurisubharmonic.
Conversely, every quasi-plurisubharmonic function $u$ on $X$ belongs to $\PSH(X,C\omega)$ for some constant $C>0$.
\end{proposition}
\Needspace{7\baselineskip}
\begin{proof}
On a coordinate ball with compact closure in a larger chart, smoothness of $\theta$ gives
\[
 \theta\leq A\ddc|z|^2
\]
for a sufficiently large $A$. Consequently,
\[
 \ddc(u+A|z|^2)\geq A\ddc|z|^2-\theta\geq0.
\]
The local distributional characterization of plurisubharmonic functions, \cite[Chapter~I, Theorem~5.8]{DemCADG}, gives a unique plurisubharmonic representative of $u+A|z|^2$. Subtracting the smooth correction gives a quasi-plurisubharmonic representative of $u$ on the ball. These representatives agree on overlaps: after a common local smooth correction, two of them become plurisubharmonic functions that are equal almost everywhere, hence everywhere. They therefore glue to a unique $\theta$-plurisubharmonic representative.

For the converse, cover $X$ by finitely many open sets $V_j\Subset U_j$, where each $U_j$ is a chart on which $u=p_j+f_j$, with $p_j$ plurisubharmonic and $f_j$ smooth. On the compact set $\overline V_j$, the smooth form $\ddc f_j$ is bounded below by $-C_j\omega$ for some constant $C_j$, so $\ddc u\geq\ddc f_j\geq-C_j\omega$ on $V_j$. Any constant $C>0$ larger than all the $C_j$ gives $\ddc u\geq-C\omega$ on $X$.
\end{proof}

For $u\in\PSH(X,\theta)$, choose coordinates $z$ and a constant $A>0$ as in \cref{prop:local-correction}, and put $p=u+A|z|^2$. Then
\[
 \theta+\ddc u=\ddc p+(\theta-A\ddc|z|^2).
\]
Thus local plurisubharmonic theory applies to $p$, while the smooth remainder $\theta-A\ddc|z|^2$ records the original background; this remainder must be kept when forming products or integrating by parts.

\begin{proposition}\label{prop:no-local-potential}\index{Background form!local potentials}
If $\theta=\ddc\rho$ on an open set for a smooth function $\rho$, then $\mathrm{d}\theta=0$ on that set. In particular, at a point where $\mathrm{d}\theta\neq0$, no neighborhood admits such a representation.
\end{proposition}
\begin{proof}
Apply $\mathrm{d}$ and use $\mathrm{d}\ddc=0$.
\end{proof}

Conversely, a closed real $(1,1)$-form has smooth local $\ddc$-potentials on sufficiently small coordinate balls, by the local $\partial\bar\partial$-lemma; this local fact does not require $X$ to admit a K\"ahler metric.

The exterior derivative also constrains the regularity of a potential whose current is pointwise proportional to a Hermitian form. In dimension at least two, differentiating such an identity expresses the derivative of the proportionality factor through smooth coefficients.

\begin{lemma}\label{lem:smooth-conformal-current}
Let $X$ be a complex manifold of dimension $n\geq2$, let $\theta$ be a smooth real $(1,1)$-form and let $\gamma$ be a Hermitian form on $X$. Suppose that real-valued $u,a\in L^1_{\mathrm{loc}}(X)$ satisfy
\[
 \theta+\ddc u=a\gamma
\]
as currents. Then $a$ and $u$ are smooth.
\end{lemma}
\begin{proof}
Taking the exterior derivative and multiplying by $\gamma^{n-2}$ gives
\[
 \mathrm{d}a\wedge\gamma^{n-1}
 =\mathrm{d}\theta\wedge\gamma^{n-2}
      -a\,\mathrm{d}\gamma\wedge\gamma^{n-2}.
\]
At each point, $\xi\mapsto\xi\wedge\gamma^{n-1}$ is an isomorphism from real one-forms to real $(2n-1)$-forms, as one sees in a unitary frame. Its inverse depends smoothly on the point. Let $B$ and $C$ be the images of $\mathrm{d}\theta\wedge\gamma^{n-2}$ and $\mathrm{d}\gamma\wedge\gamma^{n-2}$ under this inverse. They are smooth real one-forms, and
\[
 \mathrm{d}a=B-aC
\]
in the sense of distributions. Since $a\in L^1_{\mathrm{loc}}$, the right-hand side is locally integrable, so $a\in W^{1,1}_{\mathrm{loc}}$. Differentiating the equation and using the smoothness of its coefficients, we obtain inductively $a\in W^{k,1}_{\mathrm{loc}}$ for every integer $k$, and the local Sobolev embedding shows that $a$ is smooth. Consequently $\ddc u=a\gamma-\theta$ is smooth; taking its trace in local coordinates and applying elliptic regularity shows that $u$ is smooth as well.
\end{proof}

\section{Changes of background}

A positive current can be written as a smooth form plus a complex Hessian in more than one way. Changing the decomposition by a smooth $\ddc$-exact form changes the potential by a bounded function, and positive linear combinations allow us to compare potentials belonging to different backgrounds. In this section, $\eta$ denotes a second smooth real $(1,1)$-form on $X$.

\Needspace{12\baselineskip}
\begin{proposition}\label{prop:gauge}\index{Background form!change of representative}
Let $f$ be a smooth real function on $X$, and set $\theta'=\theta+\ddc f$. Then
\[
 \PSH(X,\theta)\longrightarrow\PSH(X,\theta'),
 \qquad u\longmapsto u-f
\]
is a bijection. It preserves the positive current and changes each potential by a bounded function. Moreover, $\mathrm{d}\theta'=\mathrm{d}\theta$.
\end{proposition}
\begin{proof}
For every $u\in\PSH(X,\theta)$,
\[
 \theta'+\ddc(u-f)=\theta+\ddc u.
\]
The inverse map is $v\mapsto v+f$, and $f$ is bounded because $X$ is compact. The last assertion follows from $\mathrm{d}\ddc f=0$.
\end{proof}

Thus the natural object is the class of $\theta$ modulo smooth $\ddc$-exact forms. This quotient of all smooth real $(1,1)$-forms contains nonclosed classes; the usual Bott--Chern cohomology group is the corresponding quotient of closed forms.

\Needspace{12\baselineskip}
\begin{proposition}\label{prop:convexity}\index{Plurisubharmonic function!convex combinations}
If $u\in\PSH(X,\theta)$ and $v\in\PSH(X,\eta)$, then, for $a,b\geq0$ with $a+b>0$,
\[
 au+bv\in\PSH(X,a\theta+b\eta).
\]
In particular, $\PSH(X,\theta)$ is convex. If $\eta\geq\theta$, then $\PSH(X,\theta)\subseteq\PSH(X,\eta)$.
\end{proposition}
\begin{proof}
Omitting a term whose coefficient is zero, we see that $au+bv$ is quasi-plurisubharmonic. Moreover,
\[
 a\theta+b\eta+\ddc(au+bv)
 =a(\theta+\ddc u)+b(\eta+\ddc v)\geq0.
\]
For the inclusion, add the semipositive smooth form $\eta-\theta$ to $\theta+\ddc u$.
\end{proof}

By \cref{prop:local-correction}, any two quasi-plurisubharmonic functions on $X$ belong to $\PSH(X,C\omega)$ for a common sufficiently large constant $C$. This allows us to compare their singularities within one positive background.

\section{Maxima and truncations}

Truncating a potential requires stability under maxima. For a smooth background, this follows from the local psh maximum theorem by freezing the coefficients and then letting the curvature error tend to zero.

\begin{lemma}\label{lem:max-psh}
If $u,v\in\PSH(X,\theta)$, then $\max(u,v)\in\PSH(X,\theta)$. Consequently, if $\theta\geq0$ and $u\in\PSH(X,\theta)$, then $\max(u,-t)$ is $\theta$-plurisubharmonic for every real $t$.
\end{lemma}
\begin{proof}
The maximum is quasi-plurisubharmonic by a common local smooth correction. To verify its precise lower Hessian bound, fix a positive constant $(1,1)$-form $\kappa$ in a coordinate ball. For each $\varepsilon>0$, freeze $\theta$ on smaller neighborhoods to a constant form $\theta'$ with $-\varepsilon\kappa\leq\theta-\theta'\leq\varepsilon\kappa$. Write $\theta'+\varepsilon\kappa=\ddc r$ locally. Then $u+r$ and $v+r$ are plurisubharmonic, hence so is their maximum. It follows that
\[
 \theta+2\varepsilon\kappa+\ddc\max(u,v)\geq0
\]
on each such neighborhood. Let $\varepsilon\to0$. For the last assertion, constants belong to $\PSH(X,\theta)$ when $\theta\geq0$.
\end{proof}


\begin{lemma}\label{lem:strict-set-positive-volume}
If $a,b$ are quasi-plurisubharmonic functions and $\{a<b\}$ is nonempty at a point where both are finite, then this set has positive smooth volume.
\end{lemma}
\begin{proof}
Choose a point $x$ with $a(x)<b(x)$ and a number $c$ between them. Upper semicontinuity gives $a<c$ on a neighborhood $V$ of $x$. If $\{b>c\}\cap V$ had zero volume, then $b\leq c$ almost everywhere on $V$. The canonical quasi-plurisubharmonic representative, or the submean inequality after a smooth local correction, would imply $b(x)\leq c$, a contradiction. Thus $\{a<b\}$ has positive volume.
\end{proof}

\section{Normalized compactness}

Fix a smooth positive probability volume $\mathrm{d}V$ for $L^1$ norms. The local compactness theorem for plurisubharmonic functions, recalled in \cite{SGPT}, gives an $L^1_{\loc}$-convergent subsequence or local uniform convergence to $-\infty$ for a locally uniformly upper-bounded sequence on a connected domain. A point at which the global maximum is attained excludes the second alternative, and connectedness propagates the first across $X$.

\begin{theorem}\label{thm:normalized-compactness}
For any smooth real $(1,1)$-form $\theta$, the family
\[
 \mathcal K_\theta=\{u\in\PSH(X,\theta):\sup_Xu=0\}
\]
is compact in $L^1(X,\mathrm{d}V)$. In particular, there is a constant $C_\theta$ such that
\[
 \int_X|u|\,\mathrm{d}V\leq C_\theta\quad(u\in\mathcal K_\theta).
\]
\end{theorem}
\begin{proof}
Let $(u_j)$ be a sequence in this family and choose $x_j$ with $u_j(x_j)=0$. Passing to a subsequence, suppose $x_j\to x$. Cover $X$ by finitely many connected coordinate balls, and on a slightly larger ball choose a fixed smooth correction $r_i$ for which $u_j+r_i$ is plurisubharmonic. Such corrections exist because $\theta$ is smooth. These local functions are uniformly bounded above, since $u_j\leq0$.

On a ball containing $x$ in its interior, the alternative of locally uniform convergence to $-\infty$ is impossible: the values at $x_j$ stay bounded below. Local compactness therefore gives an $L^1$-convergent subsequence on this ball. If a second ball overlaps it, local compactness on the second ball cannot yield convergence to $-\infty$ on compact subsets of the overlap, since the first limit is finite almost everywhere there. Thus a further subsequence converges in $L^1$ on the second ball. The intersection graph of a finite connected open cover of connected $X$ is connected. Repeating along that graph and taking finitely many subsequences gives local $L^1$ convergence everywhere. A partition of unity then gives convergence in $L^1(X)$ to a function $u$.

The inequalities $\ddc u_j\geq-\theta$ pass to distributions. The local canonical representatives of $u+r_i$ are plurisubharmonic and agree on overlaps after subtracting their corrections. They define a global $\theta$-plurisubharmonic representative, still denoted by $u$. Since $u_j\leq0$, it satisfies $u\leq0$.

We check that the supremum has not decreased. For a fixed small radius around $x_j$ inside the first coordinate ball, apply the submean inequality to $u_j+r_i$. Local $L^1$ convergence allows passage to the limit in the ball averages, also as their centers approach $x$: the averages of the differences tend to zero by $L^1$ convergence, and translating the indicator of a fixed ball is continuous for the integral of a fixed $L^1$ function. Letting the radius tend to zero gives
\[
 0+r_i(x)=\lim_j(u_j+r_i)(x_j)\leq(u+r_i)(x).
\]
Thus $u(x)=0$, and the limit belongs to $\mathcal K_\theta$. This proves sequential compactness, which is compactness in the metric space $L^1(X)$.

If the displayed norm bound failed, choose a sequence whose $L^1$ norms tend to infinity. Its convergent subsequence would have bounded norms, a contradiction.
\end{proof}

\begin{corollary}\label{cor:decreasing-limit}
A decreasing sequence in $\PSH(X,\theta)$ with uniformly bounded $L^1$ norms has a genuine $\theta$-plurisubharmonic limit and converges to it in $L^1$.
\end{corollary}
\begin{proof}
Local decreasing-limit theory gives a quasi-plurisubharmonic limit or the function $-\infty$ on each connected coordinate neighborhood. The $L^1$ bound excludes the latter alternative, by Fatou's lemma applied to the negative parts. The local representatives glue. Finally, the sequence is bounded above by its first member; subtracting the sequence from that member and applying monotone convergence proves $L^1$ convergence.
\end{proof}

\section{Uniform exponential integrability}

Skoda's local theorem gives an integrable exponential for each quasi-plurisubharmonic potential. Normalized compactness makes the exponent and integral bound uniform: otherwise a convergent weighted sum of counterexamples would violate the local theorem.

Write $K_\omega$ for the bound in \cref{thm:normalized-compactness}, so that normalized potentials $v\in\PSH(X,\omega)$ with $\sup_Xv=0$ satisfy
\begin{equation}\label{eq:normalized-lone}
 \int_X|v|\,\mathrm{d}V\leq K_\omega.
\end{equation}

Local Skoda integrability implies that every quasi-plurisubharmonic function $v$ on compact $X$ admits some exponent $a(v)>0$ for which $\mathrm{e}^{-a(v)v}$ is integrable. Indeed, apply the local theorem to a plurisubharmonic correction near each point, and then take the minimum exponent over a finite cover. We use the local integrability theorem recalled in \cite{SGPT}.

\begin{lemma}\label{lem:uniform-exponential}
There are constants $\alpha>0$ and $C<+\infty$, depending on $(X,\omega,\mathrm{d}V)$, such that
\[
 \int_X\mathrm{e}^{-\alpha v}\,\mathrm{d}V\leq C
 \quad\hbox{for all }v\in\PSH(X,\omega),\quad \sup_Xv=0.
\]
\end{lemma}
\begin{proof}
Suppose no such pair exists. Set $a_j=2^{-j}$ and choose normalized potentials $v_j\leq0$ with
\[
 \int_X\exp\!\left(-\frac{a_j}{j}v_j\right)\mathrm{d}V\geq2^j.
\]
The partial sums $V_N=\sum_{j\leq N}a_jv_j$ decrease, are $\omega$-plurisubharmonic, and have uniformly bounded $L^1$ norms by \eqref{eq:normalized-lone}. Their limit $V$ is therefore $\omega$-plurisubharmonic. Since all summands are nonpositive, $V\leq a_jv_j$ for every $j$.

\Needspace{6\baselineskip}
Fix $p>0$. For every $j\geq1/p$,
\[
 \int_X\mathrm{e}^{-pV}\,\mathrm{d}V
 \geq\int_X\mathrm{e}^{-pa_jv_j}\,\mathrm{d}V
 \geq\int_X\mathrm{e}^{-(a_j/j)v_j}\,\mathrm{d}V
 \geq2^j.
\]
Thus $\mathrm{e}^{-pV}$ is nonintegrable for every $p>0$, contradicting the qualitative Skoda input applied to $V$.
\end{proof}

\section{Local products and capacity}
\label{sec:psh-local-capacity}

The local Bedford--Taylor product and capacity measure the exceptional sets of plurisubharmonic functions and control bounded approximation. We recall these classical inputs from \cite{SGPT,GZ17}. Fix $1\leq k\leq n$.

\begin{theorem}[Bedford--Taylor]
\label{thm:bt-input}
If $p_1,\ldots,p_k$ are locally bounded plurisubharmonic functions on a domain in $\C^n$, their mixed product
$\ddc p_1\wedge\cdots\wedge\ddc p_k$ is a positive closed current of order zero. It is symmetric, multilinear in its current factors, and continuous under simultaneous locally bounded decreasing approximation. It does not charge pluripolar sets and is local in the plurifine topology.
\end{theorem}

Here locality means equality of the restricted currents on a plurifine open Borel set where the potentials agree. The plurifine topology is the topology generated locally by plurisubharmonic functions. It is finer than the ordinary topology and, for example, makes the strict superlevel sets of a quasi-plurisubharmonic function open.

For concentric coordinate balls $B'\Subset B$ and a Borel set $E\subset B'$, let $\operatorname{cap}(E,B)=\sup\{\int_E(\ddc v)^n:v\in\PSH(B),\ 0\leq v\leq1\}$ be the relative Bedford--Taylor capacity. We use two standard local estimates on these balls, with constants $C,a>0$ depending only on the nested balls and the fixed volume form:
\begin{enumerate}
\item If $w\leq0$ is plurisubharmonic on $B$ and $v$ is bounded and plurisubharmonic, then
\[
 \int_{B'}|w|(\ddc v)^n
 \leq C\|w\|_{L^1(B)}\|v\|_{L^\infty(B)}^n.
\]
An additive constant can first be chosen in $v$. This is the weighted Chern--Levine--Nirenberg inequality.
\item For every Borel set $E\subset B'$, the relative Bedford--Taylor capacity satisfies
\[
 \operatorname{Vol}(E)\leq C\exp\bigl(-a\operatorname{cap}(E,B)^{-1/n}\bigr),
\]
with value zero when the capacity is zero. This is the classical local volume--capacity estimate.
\end{enumerate}
These are statements of local pluripotential theory, as developed in \cite[Chapters 3 and 4]{GZ17}; the bounded product and capacity conventions also agree, up to harmless fixed constants, with \cite[Chapter 2]{SGPT}. The global Hermitian capacity estimates will be deduced in \cref{lem:global-capacity-estimates}.


Plurisubharmonic functions are quasi-continuous for the local Bedford--Taylor capacity, and Borel sets have zero capacity precisely when they are pluripolar. Uniformly bounded psh potentials converging in capacity have weakly convergent local products. The same holds after multiplying the products by uniformly bounded quasi-continuous coefficients converging in capacity. These local results, recalled in \cite{SGPT,GZ17}, also apply after a common smooth correction to quasi-plurisubharmonic functions.

\begin{theorem}\label{thm:local-mixed-input}
Let $U\subset\mathbb C^n$ be a domain and let $\beta$ be its Euclidean Hermitian form. The classical local mixed Monge--Amp\`ere inequality says that if $a,b$ are locally bounded plurisubharmonic functions and $A,B$ are nonnegative measurable densities on $U$ with
\[
  (\ddc a)^n\geq A\beta^n,
  \qquad
  (\ddc b)^n\geq B\beta^n,
\]
then, for $0\leq k\leq n$,
\[
  (\ddc a)^k\wedge(\ddc b)^{n-k}
  \geq A^{k/n}B^{(n-k)/n}\beta^n.
\]
\end{theorem}

\section{Products with smooth backgrounds}
\label{sec:psh-bounded-products}

The Bedford--Taylor construction also defines products after adding smooth forms to complex Hessians. Independence of the smooth correction makes these products global; positivity requires retaining the original lower curvature bound. Fix $1\leq k\leq n$.

\Needspace{8\baselineskip}
Let $u$ be locally bounded and $\theta$-plurisubharmonic. On a relatively compact coordinate ball with coordinates $z$, choose $A>0$ and put $q=A|z|^2$ so that $p=u+q$ is plurisubharmonic. Set $\beta=\theta-\ddc q$ and define
\begin{equation}\label{eq:bounded-expansion}
 (\theta+\ddc u)^k
 :=\sum_{j=0}^k\binom{k}{j}(\ddc p)^j\wedge\beta^{k-j}.
\end{equation}
Each summand is defined by \cref{thm:bt-input} and multiplication by a smooth form. The signed terms in this expansion should not individually be assumed positive.

\begin{proposition}\label{prop:bounded-products}
The construction \eqref{eq:bounded-expansion}, and its mixed version, is independent of the local correction $q$ and of the decomposition of each current into a smooth form and a $\ddc$-potential. It defines global products of currents with locally bounded quasi-plurisubharmonic potentials. The products are symmetric, multilinear, plurifinely local, and give no mass to pluripolar sets. If every factor is positive, their mixed product is positive.
\end{proposition}
\begin{proof}
On the overlap of two coordinate balls, the two corrected potentials differ by a smooth function. The multilinear rule
$\ddc(p+f)=\ddc p+\ddc f$, valid for bounded local products by smooth approximation, and the binomial formula identify the two expansions. The same argument handles a smooth change of the decomposition of a current. Any two potential decompositions of the same $(1,1)$-current do differ smoothly: their difference has smooth $\ddc$, and taking its local Euclidean trace gives smoothness by elliptic regularity. Symmetry, multilinearity and locality then follow term by term. Wedges of bounded Bedford--Taylor products with smooth forms have locally finite variation and do not charge pluripolar sets, so these properties also pass to the finite expansion.

It remains to prove positivity; it does not follow from the signs of the summands. Write the factors as $T_i=\theta_i+\ddc u_i$, $1\leq i\leq k$, where each $\theta_i$ is smooth and real and each $u_i$ is locally bounded and quasi-plurisubharmonic. Assume $T_i\geq0$. Fix $\varepsilon>0$ and a local positive constant form $\kappa$. After shrinking a neighborhood of a chosen point, each smooth background $\theta_i$ admits a constant, hence closed, form $\theta_i'$ such that
\[
 -\varepsilon\kappa\leq\theta_i-\theta_i'
 \leq\varepsilon\kappa.
\]
Then
$S_i=\theta_i'+\varepsilon\kappa+\ddc u_i$ is positive and closed there, with a bounded local plurisubharmonic potential. Moreover
\[
 T_i+2\varepsilon\kappa
 =S_i+(\theta_i-\theta_i'+\varepsilon\kappa).
\]
All terms on the right are positive, and the second term is smooth. Expanding the mixed product therefore gives a sum of positive currents. Although the neighborhoods depend on $\varepsilon$, they cover the original coordinate ball. The product of the $T_i+2\varepsilon\kappa$ is thus positive throughout that ball. By multilinearity it converges weakly to the product of the $T_i$ as $\varepsilon\to0$, proving positivity.
\end{proof}

The local expansion also gives continuity along uniformly locally bounded decreasing sequences, by applying \cref{thm:bt-input} to the corrected potentials and then summing the finitely many terms. Each bounded top-degree product has finite mass on compact $X$. For bounded $u\in\PSH(X,\theta)$, we write $\mathrm{MA}_\theta(u)=(\theta+\ddc u)^n$.

For $C^{1,1}$ potentials, these measure-valued products agree with the products of their measurable Hessian coefficients.

\begin{proposition}\label{prop:c11-products}
Let $u_1,\ldots,u_k$ be locally $C^{1,1}$ quasi-plurisubharmonic functions and let $\theta_1,\ldots,\theta_k$ be smooth real $(1,1)$-forms. The bounded product
\[
 (\theta_1+\ddc u_1)\wedge\cdots\wedge(\theta_k+\ddc u_k)
\]
coincides with the almost-everywhere wedge of the locally bounded measurable coefficients of its factors.
\end{proposition}
\begin{proof}
On a smaller coordinate ball, add a common smooth quadratic $q$ so that each $u_i+q$ is plurisubharmonic, and convolve with a smooth radial probability kernel. The Hessians of the convolutions are uniformly locally bounded and converge to the weak Hessians in every finite $L^s$ space. Their finite products therefore converge in $L^1$ on smaller balls. The potentials converge uniformly, so bounded Bedford--Taylor continuity identifies this limit with their pluripotential product. Expanding off the smooth correction, as in \eqref{eq:bounded-expansion}, proves the assertion for the original factors.
\end{proof}

For a locally bounded quasi-plurisubharmonic function $v$ and a smooth convex increasing function $F$ on a neighborhood of its range, the bounded Bedford--Taylor chain rule is
\begin{equation}\label{eq:bounded-convex-chain-rule}
 \ddc F(v)=F'(v)\ddc v+F''(v)\,\mathrm{d}v\wedge\dc v.
\end{equation}
The gradient term and its products with positive bounded Bedford--Taylor factors are nonnegative. These local identities follow by decreasing regularization; bounded potentials have the local Sobolev regularity needed for the gradient expression. They are part of the bounded local calculus recalled in \cite{SGPT,GZ17}. In particular, inequalities obtained from this formula pass through local smoothing with a lower-curvature error tending to zero.

\Needspace{8\baselineskip}
\section{Hermitian capacity}
\label{sec:hermitian-capacity}

Local capacity tests can be glued to the constant function on a compact manifold after a fixed rescaling. This compares local Bedford--Taylor capacity with a capacity defined by global potentials and transfers both volume estimates and normalized sublevel estimates to $X$.

For a Borel set $E\subset X$, put
\[
 \operatorname{Cap}_\omega(E)
 =\sup_{v\in\PSH(X,\omega),\ 0\leq v\leq1}
             \int_E(\omega+\ddc v)^n.
\]
The products are those of \cref{prop:bounded-products}. Their local expansion and \eqref{eq:local-cln-range} show that $\operatorname{Cap}_\omega(X)<+\infty$. On relatively compact coordinate balls, a fixed smooth correction bounds the measures of global tests by a constant times the corresponding local capacity. Conversely, the gluing construction below bounds local capacity by global capacity. Thus the quasi-continuity and capacity-zero characterizations in \cref{sec:psh-local-capacity} hold for $\operatorname{Cap}_\omega$. Convergence in this capacity means
\[
 \operatorname{Cap}_\omega(\{|u_j-u|>\delta\})\longrightarrow0
 \quad\text{for every }\delta>0,
\]
where pole sets may be omitted. The bounded-product continuity statements of \cref{sec:psh-local-capacity} pass to the Hermitian products by their local expansion and a finite cover.

For the following estimates, abbreviate $\operatorname{Cap}=\operatorname{Cap}_\omega$ and $\omega_v=\omega+\ddc v$, and let $\operatorname{Vol}$ denote the volume for the fixed smooth measure $\mathrm{d}V$.

\begin{lemma}\label{lem:global-capacity-estimates}
There are constants $a,C>0$ such that
\begin{equation}\label{eq:global-volume-capacity}
 \operatorname{Vol}(E)\leq C\exp\bigl(-a\operatorname{Cap}(E)^{-1/n}\bigr)
\end{equation}
for every Borel $E\subset X$. Also, for $u\in\PSH(X,\omega)$ with $\sup_Xu=0$ and $t>1$,
\begin{equation}\label{eq:global-capacity-tail}
 \operatorname{Cap}(\{u<-t\})\leq C/t.
\end{equation}
\end{lemma}
\begin{proof}
Take a finite cover by smaller coordinate balls $B'_j\Subset B_j$. A local capacity test $0\leq v\leq1$ on $B_j$ can be made into a global test as follows. Choose a smooth function $\chi_j$ equal to zero near $\overline{B'_j}$, equal to $-2$ near $\partial B_j$, and with $-2\leq\chi_j\leq0$. Choose $\delta>0$ small enough that
\[
 \omega+\delta\ddc\chi_j\geq0,\qquad 3\delta/2\leq1.
\]
The function
\[
 V=\max\{\delta(v+\chi_j),-\delta/2\}+\delta/2
 \quad\hbox{on }B_j,
 \qquad V=0\quad\hbox{outside }B_j
\]
is globally $\omega$-plurisubharmonic and lies between zero and one. Near the boundary the constant branch is selected. On $B'_j$ we have $\omega_V\geq\delta\ddc v$. Positivity of bounded products therefore gives
\[
 \operatorname{cap}(E\cap B'_j,B_j)
 \leq\delta^{-n}\operatorname{Cap}(E).
\]
The local volume estimate, followed by summation over the finite cover, proves \eqref{eq:global-volume-capacity}.

For the second estimate, choose on each $B_j$ a smooth negative quadratic function $r_j$ with $\ddc r_j\geq\omega$. For every global test $0\leq v\leq1$, the functions $u+r_j\leq0$ and $v+r_j$ are plurisubharmonic. Thus
\[
 \begin{split}
 \int_{\{u<-t\}}\omega_v^n
 &\leq t^{-1}\int_X(-u)\omega_v^n\\
 &\leq t^{-1}\sum_j\int_{B'_j}(-u-r_j)
                    \bigl(\ddc(v+r_j)\bigr)^n\\
 &\leq \frac Ct\left(1+\int_X|u|\,\mathrm{d}V\right).
 \end{split}
\]
Here the second inequality uses $\omega_v\leq\ddc(v+r_j)$, and the third is the local weighted estimate; the local test functions have a common bounded range. The normalized $L^1$ bound in \cref{thm:normalized-compactness} finishes the proof after taking the supremum over $v$.
\end{proof}

\begin{corollary}\label{cor:lp-capacity-square}
Fix $p>1$. There is $C_p$ such that every nonnegative $F\in L^p(X,\mathrm{d}V)$ satisfies
\[
 \int_E F\,\mathrm{d}V\leq C_p\|F\|_p\operatorname{Cap}(E)^2
\]
for every Borel set $E$.
\end{corollary}
\begin{proof}
H\"older's inequality and \eqref{eq:global-volume-capacity} bound the integral by
\[
 C\|F\|_p\exp(-a\operatorname{Cap}(E)^{-1/n}/p'),
 \qquad p'=p/(p-1).
\]
The quotient of this exponential by $x^2$ is bounded on $0<x\leq\operatorname{Cap}(X)<\infty$. Capacity-zero sets have zero smooth volume.
\end{proof}

The same proof bounds an $L^p$ density by any fixed positive power of capacity. The choice of the exponent two will be convenient for the sublevel iteration in \cref{lem:backwards-capacity}.

For quasi-continuous obstacles, a decreasing lower semicontinuous approximation preserves the original bounds and confines the error to sets of vanishing capacity.

\begin{lemma}
\label{lem:quasicont-lsc-approximation}
Every bounded quasi-continuous function $f$ admits bounded lower semicontinuous functions $f_j$ with common bounds such that $f_j\downarrow f$ quasi-everywhere and $f_j\to f$ in capacity.
\end{lemma}
\begin{proof}
Choose open sets $O_j$ with capacity at most $2^{-j}$ such that $f$ is continuous on $X\setminus O_j$, and set $G_j=\bigcup_{k\geq j}O_k$. Then $G_j$ decreases, its capacity tends to zero, and $f$ is continuous on its closed complement. Choose $M\geq\sup_Xf$ and define
\[
 f_j=f\ \hbox{on }X\setminus G_j,\qquad f_j=M\ \hbox{on }G_j.
\]
At a boundary point of $G_j$, limits from the closed complement give $f$, while values from inside are at least that value. Thus $f_j$ is lower semicontinuous. The functions decrease to $f$ outside $\bigcap_jG_j$, a capacity-zero set. Their difference from $f$ is supported in $G_j$, which also proves convergence in capacity.
\end{proof}

\Needspace{10\baselineskip}
\section{Comparing local regularizations}

To glue local convolutions, we need them to agree up to an error tending uniformly to zero on overlaps. A bounded plurisubharmonic function need not be continuous, so this comparison cannot come from uniform convergence to the original function. Instead, logarithmic convexity of its suprema on balls controls the change of scale and the change of coordinates.

\Needspace{9\baselineskip}
Let $U\subset\mathbb C^n$ be a domain and $u$ a bounded plurisubharmonic function on $U$. On a fixed relatively compact subset, for small $\delta>0$ put
\[
 S_\delta u(z)=\sup_{|\zeta-z|\leq\delta}u(\zeta),\qquad
 A_\delta u(z)=\fint_{|\zeta-z|=\delta}u(\zeta)\,\mathrm{d}\sigma(\zeta),
 \qquad u_\delta=u*\rho_\delta,
\]
where $\mathrm{d}\sigma$ is Euclidean surface measure and $\rho$ is a smooth nonnegative radial mollifier of integral one, supported in the unit ball, with $\rho_\delta(\zeta)=\delta^{-2n}\rho(\zeta/\delta)$. A fixed larger neighborhood contains all balls below, and $B$ denotes a bound for the oscillation of $u$ there.

\Needspace{16\baselineskip}
\begin{lemma}\label{lem:bounded-local-regularization}
For each fixed $A>1$,
\[
 S_{A\delta}u-S_\delta u\longrightarrow0,
 \qquad S_\delta u-u_\delta\longrightarrow0
\]
uniformly on compact subsets as $\delta\downarrow0$.
\end{lemma}
\begin{proof}
The maximum of a plurisubharmonic function on concentric balls is convex as a function of the logarithm of the radius. This follows by restricting to each complex line through the center and comparing on an annulus with the harmonic function affine in $\log|\zeta|$. For a fixed admissible radius $r>A\delta$, convexity gives
\begin{equation}\label{eq:bounded-log-radius}
 0\leq S_{A\delta}u-S_\delta u
 \leq\frac{B\log A}{\log(r/\delta)}.
\end{equation}
This proves the first assertion.

We compare the maximum with the spherical mean using the Poisson kernel in real dimension $2n$. For a point at distance at most $\delta/2$ from the center, the normalized Poisson kernel of the ball of radius $\delta$ is bounded below by a positive dimensional constant $c_n$. Since $u\leq S_\delta u$, the submean comparison with its harmonic extension gives
\[
 S_\delta u-S_{\delta/2}u
 \geq c_n(S_\delta u-A_\delta u).
\]
These inequalities are valid for bounded plurisubharmonic functions by local smoothing, followed by a decreasing limit. Equation \eqref{eq:bounded-log-radius} shows that $S_\delta u-A_\delta u$ tends uniformly to zero.

For $0<t\leq1$, the spherical mean is increasing with radius, and
\[
 0\leq A_\delta u-A_{t\delta}u
 \leq (S_\delta u-S_{t\delta}u)
                +(S_{t\delta}u-A_{t\delta}u).
\]
The first term is at most
$B\log(1/t)/\log(r/(t\delta))$ by logarithmic convexity of $S$, and the second is at most $C_nB/\log(r/(t\delta))$ by the preceding estimate. Write the radial convolution as an average of $A_{t\delta}u$ with the probability density
$\widehat\rho(t)=|\mathbb S^{2n-1}|t^{2n-1}\rho(t)$. Since
$\int_0^1(1+\log(1/t))\widehat\rho(t)\,\mathrm{d}t<\infty$, integration gives
\[
 0\leq A_\delta u-u_\delta
 \leq\frac{C_{n,\rho}B}{\log(r/\delta)}.
\]
Combining the two comparisons proves the second assertion.
\end{proof}

\Needspace{9\baselineskip}
\begin{corollary}\label{cor:coordinate-convolution}
Let $U,V\subset\mathbb C^n$ be domains and let $F:U\to V$ be biholomorphic and $u$ bounded and plurisubharmonic on $U$. The difference
\[
 u_\delta-\bigl((u\circ F^{-1})*\rho_\delta\bigr)\circ F
\]
tends locally uniformly to zero as $\delta\downarrow0$.
\end{corollary}
\begin{proof}
On a compact subset, $F$ and its inverse have uniformly bounded first derivatives. Thus, for a fixed $A>1$, the inverse image of a radius-$\delta$ ball about $F(z)$ lies between balls about $z$ of radii $\delta/A$ and $A\delta$. The two supremum regularizations are therefore bounded by one another at these enlarged radii. The first assertion of \cref{lem:bounded-local-regularization} makes their difference uniformly small. Its second assertion compares each of the two convolutions with the corresponding supremum regularization.
\end{proof}

\Needspace{12\baselineskip}
\section{Gluing with a small curvature error}

The overlap estimate permits gluing by a smooth maximum. Small boundary corrections make each local branch irrelevant near the edge of its chart, while convexity of the maximum preserves a common lower curvature bound. Use $\omega$ to measure the curvature error.

For an integer $N\geq1$, choose independent real random variables $Z_1,\ldots,Z_N$ with a common smooth even probability density supported in $[-1,1]$, and define
\[
 M_\tau(t_1,\ldots,t_N)
     =\mathbb E\max_i(t_i+\tau Z_i)\quad(\tau>0).
\]
Convolution makes $M_\tau$ smooth. It is convex, increasing in every argument, and, with $\mathbf1=(1,\ldots,1)$, satisfies
\[
 \max_i t_i\leq M_\tau(t)\leq\max_i t_i+\tau,
 \qquad M_\tau(t+c\mathbf1)=M_\tau(t)+c
\]
for $t\in\mathbb R^N$ and $c\in\mathbb R$. If $t_i<\max_{j\ne i}t_j-2\tau$, the $i$th argument can be removed without changing the value. Also $M_\tau(t)$ is nondecreasing in $\tau$: it is a convex even function of $\tau$ because the distribution of $Z$ is symmetric.

The second identity expresses translation invariance. If smooth real functions $f_i$, $1\leq i\leq N$, satisfy $\ddc f_i\geq-\gamma$ for a common smooth real $(1,1)$-form $\gamma$, then
\begin{equation}\label{eq:smooth-max-curvature}
 \ddc M_\tau(f_1,\ldots,f_N)\geq-\gamma.
\end{equation}
Indeed, the first-derivative coefficients are nonnegative and sum to one, by the translation identity, while the real Hessian of the convex function $M_\tau$ contributes a semipositive $(1,1)$-form in the chain rule.

\Needspace{8\baselineskip}
\begin{proposition}\label{prop:bounded-almost-regularization}
Let $\theta$ be smooth and real, $\omega$ Hermitian, and
$u\in\PSH(X,\theta)\cap L^\infty(X)$. There are smooth functions $v_j$ and positive numbers $\varepsilon_j\downarrow0$ such that
\[
 v_j\downarrow u,\qquad
 \theta+\varepsilon_j\omega+\ddc v_j\geq0.
\]
\end{proposition}
\Needspace{9\baselineskip}
\begin{proof}
Fix $\epsilon>0$. By continuity of $\theta$, choose finitely many nested coordinate neighborhoods
$V_i\Subset W_i\Subset U_i$, with the $V_i$ covering $X$, and smooth quadratic functions $r_i$ such that
\[
 0\leq\ddc r_i-\theta\leq\epsilon\omega\quad\hbox{on }U_i.
\]
To make this choice, freeze the coefficients of $\theta$ at the center, add a small positive constant form, and shrink the coordinate neighborhood. Put $p_i=u+r_i$, a bounded plurisubharmonic function, and let $p_{i,\delta}$ be its convolution in $U_i$.

On an overlap, express both $p_i$ and $p_j$ in the $j$th coordinates. Their difference is the smooth function $r_i-r_j$, so the difference of their convolutions tends uniformly to $r_i-r_j$. By \cref{cor:coordinate-convolution}, changing the coordinates used to convolve $p_i$ introduces another uniformly vanishing error. Thus, on overlaps of the closures of the $W_i$,
\begin{equation}\label{eq:overlap-regularization}
 (p_{i,\delta}-r_i)-(p_{j,\delta}-r_j)\longrightarrow0
 \quad\hbox{uniformly}.
\end{equation}
Choose $\eta_i\in C^\infty(U_i,[-1,0])$, zero on $V_i$ and equal to $-1$ near the boundary of $W_i$, and $C\geq1$ with $\ddc\eta_i\geq-C\omega$. The cover, the functions $r_i,\eta_i$ and the constant $C$ may depend on $\epsilon$; they are fixed while $\delta$ tends to zero. Set
\[
 f_{i,\delta}=p_{i,\delta}-r_i+(\epsilon/C)\eta_i
 \quad\hbox{on }W_i.
\]
Each satisfies $\ddc f_{i,\delta}\geq-\theta-2\epsilon\omega$. Near a boundary point of $W_i$, some $V_j$ contains that point. Equation \eqref{eq:overlap-regularization} implies, for all small $\delta$,
\[
 f_{i,\delta}<f_{j,\delta}-\epsilon/(2C)
\]
on a neighborhood of that boundary point. Finiteness of the cover makes this separation uniform.

Take $\tau(\delta)\downarrow0$ as $\delta\downarrow0$, with $\tau(\delta)<\epsilon/(8C)$, and define on each point the smooth maximum of the branches whose $W_i$ contain it. A branch disappears only near the boundary of its $W_i$, where it is below another branch by more than $2\tau$. The removal property of $M_\tau$ therefore proves that these formulas define a global smooth function $\psi_{\epsilon,\delta}$. Equation \eqref{eq:smooth-max-curvature} gives
\[
 \theta+2\epsilon\omega+\ddc\psi_{\epsilon,\delta}\geq0.
\]
At every point there is a branch with $\eta_i=0$, and $p_{i,\delta}\geq p_i$. Hence $\psi_{\epsilon,\delta}\geq u$. For fixed $\epsilon$, the local convolutions decrease with $\delta$, and so does the smooth maximum when both its arguments and $\tau$ decrease. Its pointwise limit is
\[
 \max_{i:x\in W_i}\bigl(u(x)+(\epsilon/C)\eta_i(x)\bigr)=u(x).
\]
Thus $\psi_{\epsilon,\delta}\downarrow u$ as $\delta\downarrow0$, also in $L^1$.

Choose a decreasing diagonal as follows. Take $\epsilon_j\downarrow0$ and put
$v_j=\psi_{\epsilon_j,\delta_j}+2^{-j}$. Choose $\delta_j$ so small that
\[
 \|\psi_{\epsilon_j,\delta_j}-u\|_{L^1}\leq2^{-j},
 \qquad v_j<v_{j-1}\quad(j\geq2).
\]
The second choice is possible because $v_{j-1}\geq u+2^{-j+1}$. More explicitly, the continuous functions
$\psi_{\epsilon_j,\delta}+2^{-j}-v_{j-1}$ decrease pointwise to a function bounded above by $-2^{-j}$. Their nonnegative level sets are decreasing compact sets with empty intersection, so one of them is empty. The resulting sequence decreases, converges to $u$ in $L^1$, and stays above $u$. Its pointwise limit equals $u$ by uniqueness of canonical quasi-plurisubharmonic representatives. Taking $\varepsilon_j=2\epsilon_j$ proves the proposition.
\end{proof}

\Needspace{16\baselineskip}
\section{Regularization in a fixed Hermitian background}

For $\theta=\omega$, rescaling the approximants absorbs the curvature error into the positive background. Subtracting a constant first makes the approximants negative, so that this rescaling preserves their decreasing order.

\begin{theorem}\label{thm:bounded-exact-regularization}\index{Regularization!bounded potentials}
For $u\in\PSH(X,\omega)\cap L^\infty(X)$ there are smooth functions $u_j\downarrow u$ such that
\[
 \omega+\ddc u_j>0.
\]
The functions $u_j$ are uniformly bounded.
\end{theorem}
\begin{proof}
Subtract a constant so that $u\leq-2$. Apply \cref{prop:bounded-almost-regularization} with $\theta=\omega$, obtaining smooth $v_j\downarrow u$ and positive $\varepsilon_j\downarrow0$ with $(1+\varepsilon_j)\omega+\ddc v_j\geq0$. The compact sets $\{v_j\geq-1\}$ decrease to the empty set, since $u\leq-2$. Discarding finitely many terms therefore gives $v_j<0$ everywhere. Write
\[
 u_j=\frac{v_j}{1+2\varepsilon_j}.
\]
Then
\[
 \omega+\ddc u_j
 \geq\frac{\varepsilon_j}{1+2\varepsilon_j}\omega>0.
\]
Because $v_j$ is negative and decreases while the positive denominator decreases to one, $u_j$ also decreases and converges to $u$. It lies between $u$ and its first smooth member, which gives uniform boundedness. Restore the subtracted constant.
\end{proof}

This regularization supplies the approximation used for products and mass in \cref{chap:products}. For an arbitrary smooth $\theta$, \cref{prop:bounded-almost-regularization} retains a vanishing curvature error.


\section{Singularity type and multiplier ideals}\label{sec:psh-singularity-types}

There are two natural ways to compare the singularities of quasi-plurisubharmonic functions on $X$. The first asks whether one function is bounded above by the other up to an additive constant; it leads to the ordinary singularity type. The second records only which holomorphic functions are locally square integrable against the weights $\mathrm{e}^{-ku}$, $k>0$; it leads to multiplier ideals. Domination implies the corresponding inclusion of multiplier ideals. A logarithmic pole on a disk shows that a single weight loses information, which is why the second comparison uses all positive weights.

\begin{definition}\label{def:singularity-order}\index{Singularity type!ordinary}\index{Singularity order}
Let $u,v$ be quasi-plurisubharmonic functions on $X$. We write $u\preceq v$ if there is a constant $C\in\R$ such that $u\leq v+C$ on $X$, and $u\sim v$ if both $u\preceq v$ and $v\preceq u$; in the latter case $u$ and $v$ have the same \emph{singularity type}. When this has to be distinguished from the multiplier-ideal type introduced below, we speak of the \emph{ordinary singularity type}, or simply the \emph{ordinary type}.
\end{definition}

Thus $u\preceq v$ means that $u$ is at least as singular as $v$. For instance, a quasi-plurisubharmonic function with a logarithmic pole lies below every bounded potential in this order, but not conversely.

\begin{proposition}\label{prop:singularity-elementary}
The relation $\preceq$ is a preorder, and $\sim$ is an equivalence relation. If $u$ is quasi-plurisubharmonic and $f$ is a smooth function on $X$, then $u+f\sim u$. In particular, all smooth functions on $X$ have the same singularity type.
\end{proposition}
\begin{proof}
Reflexivity follows with $C=0$. If $u\leq v+C$ and $v\leq w+D$, then $u\leq w+C+D$, which proves transitivity. As for every preorder, comparison in both directions defines an equivalence relation. Since $X$ is compact, a smooth function $f$ satisfies $|f|\leq A$ for some constant $A$, and then $u-A\leq u+f\leq u+A$. Taking $u=0$ shows that every smooth function has the singularity type of the constants.
\end{proof}

Some later results require the singularities of a potential to be concentrated on a small closed set. A closed subset $A\subseteq X$ is \emph{complete pluripolar} if every point of $X$ has a neighborhood $U$ with a plurisubharmonic function $p$ on $U$ such that $A\cap U=\{p=-\infty\}$. A quasi-plurisubharmonic function has \emph{small unbounded locus} if it is locally bounded outside a closed complete pluripolar set.\index{Small unbounded locus}

\begin{definition}\label{def:multiplier-ideal}\index{Singularity type!multiplier ideal}\index{Multiplier ideal!definition}
Let $\mathcal O_X$ denote the sheaf of holomorphic functions on $X$. For a quasi-plurisubharmonic function $u$ on $X$, the \emph{multiplier ideal} $\mathcal I(u)\subseteq\mathcal O_X$ is the ideal sheaf whose sections over an open set $U$ are the holomorphic functions $g$ on $U$ for which $|g|^2\mathrm{e}^{-u}$ is locally integrable on $U$ with respect to a smooth positive volume form. Its stalk at $x$ is denoted by $\mathcal I(u)_x$.
\end{definition}

Changing the volume form does not change local integrability: on relatively compact coordinate neighborhoods, the ratio of two smooth positive volume forms is bounded above and below by positive constants. In particular, the definition does not involve a choice of Hermitian metric, and it applies verbatim to quasi-plurisubharmonic functions on any complex manifold.

\begin{proposition}\label{prop:ideal-order}
If $u\preceq v$, then $\mathcal I(ku)\subseteq\mathcal I(kv)$ for every real $k>0$. If $u\sim v$, then these ideals agree for every $k>0$.
\end{proposition}
\begin{proof}
Let $g$ be a local holomorphic function on $X$. If $u\leq v+C$, then $\mathrm{e}^{-kv}\leq \mathrm{e}^{kC}\mathrm{e}^{-ku}$. Multiplying by $|g|^2$ shows that local integrability of $|g|^2\mathrm{e}^{-ku}$ implies local integrability of $|g|^2\mathrm{e}^{-kv}$. If $u\sim v$, the same argument applies in both directions.
\end{proof}

\begin{example}\label{ex:ideal-disc}\index{Multiplier ideal!logarithmic pole}
On a disk in $\mathbb C$ centered at the origin, with coordinate $z$, let $\lambda\geq0$ and put $u(z)=\lambda\log|z|^2$ for $\lambda>0$, while $u=0$ for $\lambda=0$. Then
\[
 \mathcal I(u)_0=(z^{\lfloor\lambda\rfloor}).
\]
\Needspace{8\baselineskip}
Indeed, if a nonzero holomorphic germ has order $m$ at the origin, its squared modulus is comparable to $|z|^{2m}$ near the origin. In polar coordinates, its weighted integral over a disk of sufficiently small radius $\epsilon>0$ is finite precisely when
\[
 \int_0^\epsilon r^{2m-2\lambda+1}\,\mathrm{d}r<\infty,
 \qquad\text{equivalently}\qquad m>\lambda-1.
\]
The smallest nonnegative integer satisfying this strict inequality is $\lfloor\lambda\rfloor$. When $\lambda$ is a positive integer, the borderline exponent $-1$ occurs for $m=\lambda-1$ and gives a divergent integral; this is where the strictness of the inequality matters.
\end{example}

The example also shows that a single multiplier ideal does not determine the singularity: different coefficients $\lambda$ can have the same integer part. For a real weight $k>0$, the same calculation applied to $ku$ gives $\mathcal I(ku)_0=(z^{\lfloor k\lambda\rfloor})$. The ideals at all weights $k>0$ recover $\lambda$: if $0\leq\lambda<\mu$, the integer parts of $k\lambda$ and $k\mu$ differ as soon as $k(\mu-\lambda)>1$.

We therefore compare multiplier ideals at all positive weights. For quasi-plurisubharmonic functions $u,v$, we write $u\preceq_{\mathcal I}v$ if $\mathcal I(ku)\subseteq\mathcal I(kv)$ for every real $k>0$, and $u\sim_{\mathcal I}v$ if $\mathcal I(ku)=\mathcal I(kv)$ for every real $k>0$.\index{Multiplier ideal!ideal order and ideal type} We call $\preceq_{\mathcal I}$ the \emph{ideal order}; if $u\sim_{\mathcal I}v$, then $u$ and $v$ have the same \emph{multiplier-ideal type}, or \emph{ideal type}. By \cref{prop:ideal-order}, $u\preceq v$ implies $u\preceq_{\mathcal I}v$, and $u\sim v$ implies $u\sim_{\mathcal I}v$. These relations make sense on any complex manifold. Their relation to envelopes and to global masses requires further arguments; see \cref{chap:hermitian-igood}.

\section{Analytic types and restriction}\label{sec:analytic-types}\index{Lelong number}\index{Analytic singularity}

The simplest singularities of quasi-plurisubharmonic functions are logarithmic poles along the common zeros of finitely many holomorphic functions. The logarithmic growth of a general singularity at a point is measured by its Lelong number. Let $u$ be a quasi-plurisubharmonic function on a complex manifold and let $x$ be a point. Choose holomorphic coordinates centered at $x$ and a smooth function $q$ such that $p=u+q$ is plurisubharmonic near $x$, and write $B(x,r)$ for the Euclidean coordinate ball of radius $r$. The \emph{Lelong number} of $u$ at $x$ is
\[
 \nu(u,x)=\lim_{r\downarrow0}
       \frac{\sup_{B(x,r)}p}{\log r^2}.
\]
By the local growth theorem recalled in \cite{SGPT}, the limit exists and does not depend on the coordinates or on the correction $q$. With this normalization, $\log|z|^2$ has Lelong number one at the origin. For an irreducible analytic subset $Z$, the \emph{generic Lelong number} of $u$ along $Z$ is
\[
 \nu(u,Z)=\inf_{x\in Z}\nu(u,x).
\]
By Siu's theorem, recalled in \cref{subsec:psh-lelong-levels}, $\nu(u,x)=\nu(u,Z)$ for all $x\in Z$ outside a countable union of proper analytic subsets of $Z$.

For analytic singularities we follow the convention of \cite{SGPT}: we say that a quasi-plurisubharmonic function $u$ has \emph{analytic singularities} if near every point there are holomorphic functions $f_1,\ldots,f_N$, not all identically zero, and a rational number $c>0$ such that
\[
 u=c\log\bigl(|f_1|^2+\cdots+|f_N|^2\bigr)+O(1),
 \qquad c\in\mathbb Q_{>0}.
\]
The remainder is only required to be bounded; if it can be chosen smooth, we say that $u$ has analytic singularities \emph{with smooth remainder}. Given a coherent ideal sheaf $\mathfrak a$ and a common coefficient $c$, we say that $u$ has \emph{analytic type} $(c,\mathfrak a)$ if this holds with local generators $f_1,\ldots,f_N$ of $\mathfrak a$. The choice of generators does not matter: two finite systems of generators of the same ideal change the logarithm only by a bounded amount on smaller relatively compact open sets.

\begin{proposition}\label{prop:analytic-type-calculus}
Analytic singularities are preserved by addition, finite maxima, and multiplication by positive rational numbers. Let $u$ have analytic type $(c,\mathfrak a)$, and let $Y\subset X$ be a connected complex submanifold on which $\mathfrak a$ does not vanish identically. Then the restriction $u|_Y$ is a quasi-plurisubharmonic function on $Y$, in particular not identically $-\infty$, and it has analytic type
\[
 (c,\mathfrak a\mathcal O_Y).
\]
Here $\mathfrak a\mathcal O_Y$ is the image ideal, generated by the restrictions to $Y$ of local sections of $\mathfrak a$.
\end{proposition}
\begin{proof}
For multiplication by a positive rational number, multiply the coefficient. Consider two presentations with coefficients $c_1,c_2$ and ideals $\mathfrak a_1,\mathfrak a_2$, and choose a positive rational number $d$ with $c_1=k d$ and $c_2=\ell d$, where $k,\ell$ are positive integers. Raising a sum of squared moduli of generators of $\mathfrak a_1$ to the power $k$ gives, up to bounded factors, the corresponding sum for $\mathfrak a_1^k$; similarly for $\mathfrak a_2^\ell$. Products of generators of $\mathfrak a_1^k$ and $\mathfrak a_2^\ell$ therefore present the sum with coefficient $d$, since a product of sums of squares is a sum of squares of pairwise products.

\Needspace{7\baselineskip}
For the maximum, apply the elementary inequalities
\[
 \max\{A,B\}\leq A+B\leq2\max\{A,B\},\qquad A,B\geq0,
\]
to the two sums of squared moduli of generators of $\mathfrak a_1^k$ and $\mathfrak a_2^\ell$. Taking logarithms shows that the maximum differs from $d$ times the logarithm of their total by a bounded amount. These are local assertions; on the compact manifold, finitely many charts give the corresponding comparisons of ordinary singularity types.

For the restriction, restrict to $Y$ the two-sided bounds that define the remainder. Since $Y$ is connected and is not contained in the zero set of $\mathfrak a$, the common zero set of the restricted generators is a proper analytic subset of $Y$, so $u|_Y$ is finite on a dense open subset of $Y$. After a smooth local correction, the restriction theorem for plurisubharmonic functions shows that $u|_Y$ is quasi-plurisubharmonic, and the restricted bounds give the asserted presentation.
\end{proof}

For an arbitrary quasi-plurisubharmonic function, a vanishing generic Lelong number along a submanifold does not prevent the restriction from being identically $-\infty$. For example, on a small polydisk in $\mathbb C^n$ with coordinates $z=(z_1,\ldots,z_n)$ and $|z_1|<\mathrm{e}^{-1}$, the function
\[
 u(z)=-\sqrt{-\log|z_1|^2}
\]
is plurisubharmonic, because $t\mapsto-\sqrt{-t}$ is increasing and convex for $t<0$. Its Lelong number vanishes at every point of $\{z_1=0\}$, but its restriction to this hyperplane is identically $-\infty$. For the analytic approximants constructed below, the analytic hypothesis in \cref{prop:analytic-type-calculus} excludes this behavior.

\subsection{Lelong level sets}
\label{subsec:psh-lelong-levels}\index{Siu analyticity theorem}

Siu's analyticity theorem states that, for every quasi-plurisubharmonic function $u$ and every $c>0$, the upper level set $\{x:\nu(u,x)\geq c\}$ is an analytic subset. After a smooth local correction, this is the theorem for closed positive currents, which we use in the form recalled in \cite{SGPT}; see \cref{sec:history-psh} for the original source.

For a closed positive current, the Lelong monotonicity formula also bounds the trace mass of small balls. The resulting Hausdorff estimate shows that a closed positive $(j,j)$-current puts no mass on sets of zero Hausdorff $2(n-j)$-measure.

\Needspace{10\baselineskip}
\begin{corollary}\label{cor:closed-current-hausdorff-null}\index{Closed positive current!Hausdorff support}
Let $T$ be a closed positive $(j,j)$-current on a complex manifold of dimension $n$, with $1\leq j<n$. Its coefficient measures vanish on every Borel set of zero Hausdorff $2(n-j)$-measure.
\end{corollary}
\begin{proof}
In a coordinate chart put $\kappa=\ddc|z|^2$ and $\mu=T\wedge\kappa^{n-j}$. By the Lelong monotonicity formula in bidimension $(n-j,n-j)$, the function $r\mapsto r^{-2(n-j)}\mu(B(x,r))$ is nondecreasing; see \cite[Chapter III, formulas (5.5)--(5.8), pp.~158--159]{DemCADG}. Fix a relatively compact inner chart and a radius $r_0>0$ such that all balls of radius $r_0$ centered in it lie in a fixed compact subset of the chart. Their masses are bounded by the mass of that compact set. Hence
\[
 \mu(B(x,r))\leq Cr^{2(n-j)},\qquad 0<r\leq r_0,
\]
uniformly for centers in the inner chart. A Borel set in the inner chart with zero Hausdorff $2(n-j)$-measure can be covered by balls $B(x_i,r_i)$ with $\sum_ir_i^{2(n-j)}$ arbitrarily small, so it has zero $\mu$-mass. Since $T$ is positive, the total variation of each of its coefficient measures is bounded by a constant multiple of $\mu$. A countable cover by charts proves the assertion.
\end{proof}

The non-strict threshold in Siu's theorem is essential, even on a compact K\"ahler manifold.

\begin{example}\label{ex:strict-lelong-level}\index{Lelong number!strict level sets}
Let $X=\mathbb P^2$, with homogeneous coordinates $[x:y:z]$, and let $\sigma$ be the Fubini--Study form normalized by $\int_X\sigma^2=1$. Put $\|Z\|^2=|x|^2+|y|^2+|z|^2$, $a_j=2^{-j}$ and $\lambda_j=2^{-j-1}$ for $j\geq1$. The global function
\[
 u([Z])=\log\frac{|y|^2}{\|Z\|^2}
       +\sum_{j\geq1}\lambda_j
          \log\frac{|x-a_jz|^2}{(1+a_j^2)\|Z\|^2}
\]
belongs to $\PSH(X,\frac32\sigma)$. For the line $D=\{y=0\}$ and the points $p_j=[a_j:0:1]$, it satisfies
\[
 \{\nu(u,\cdot)\geq1\}=D,
 \qquad \{\nu(u,\cdot)>1\}=\{p_j:j\geq1\}.
\]
The second set is not analytic.
\end{example}
\begin{proof}
For a unit linear form $\ell$, the function $v_\ell=\log(|\ell(Z)|^2/\|Z\|^2)$ is nonpositive and satisfies $\ddc v_\ell=[\ell=0]-\sigma$. Its $L^1(X,\sigma^2)$ norm is finite and independent of $\ell$: the group $U(3)$ acts on homogeneous coordinates, preserves $\sigma$, and acts transitively on unit linear forms. Since $\sum_j\lambda_j=1/2$, the displayed partial sums decrease to an $L^1$ function. Their common curvature lower bound is $-\frac32\sigma$, so their limit is quasi-plurisubharmonic and
\[
 T:=\tfrac32\sigma+\ddc u=[D]+\sum_{j\geq1}\lambda_j[L_j],
 \qquad L_j=\{x-a_jz=0\}.
\]
The series of positive closed currents converges: every line has $\sigma$-trace mass one and the coefficients are summable.

We compute the Lelong numbers from the local mass formula. In affine coordinates $w=(w_1,w_2)$ centered at a point $q$, put $\kappa=\ddc|w|^2$. An affine complex line $H$ at distance $d$ from $q$ satisfies
\[
 \int_{H\cap B(q,r)}\kappa=(r^2-d^2)_+.
\]
Consequently $r^{-2}\int_{H\cap B(q,r)}\kappa$ is at most one and tends to $\mathbf1_H(q)$ as $r\downarrow0$; a line outside the chart contributes zero. Tonelli's theorem and dominated convergence for the summable coefficients therefore give
\[
 \nu(u,q)=\nu(T,q)
       =\mathbf1_D(q)+\sum_{j\geq1}\lambda_j\mathbf1_{L_j}(q).
\]
Here a smooth local correction does not change the Lelong number. This computation includes points where infinitely many lines accumulate.

The lines $L_j$ meet $D$ at the distinct points $p_j$ and meet one another only at $p_\infty=[0:1:0]\notin D$. Thus $\nu(u,p_j)=1+\lambda_j$, whereas $\nu(u,q)=1$ for $q\in D\setminus\{p_j:j\geq1\}$. Outside $D$, the Lelong number is at most $\sum_j\lambda_j=1/2$. The two level-set identities follow. Finally, $p_j\to p_*=[0:0:1]$, and $\nu(u,p_*)=1$. The strict level set is not closed and hence is not analytic.
\end{proof}


\section{The local Bergman estimates}

Demailly's approximation of plurisubharmonic functions by functions with analytic singularities is built from Bergman kernels of weighted spaces of holomorphic functions. In this section we recall the local construction and its estimates. Fix a bounded pseudoconvex domain $\Omega\subset\mathbb C^n$ and a plurisubharmonic function $\psi$ on $\Omega$. Write $\mathrm{d}\lambda$ for Euclidean volume and $B(x,r)$ for the Euclidean ball of radius $r>0$ about $x$. For a positive integer $m$, define the weighted Hilbert space of holomorphic functions
\[
 \mathcal H_m=\left\{f\in\mathcal O(\Omega):
          \|f\|_m^2:=\int_\Omega|f|^2\mathrm{e}^{-m\psi}\,\mathrm{d}\lambda<\infty\right\}.
\]
\Needspace{8\baselineskip}
Choose an orthonormal basis $(g_{m,j})_j$ of $\mathcal H_m$ and define $\beta_m=m^{-1}\log\sum_j|g_{m,j}|^2$. The sum is the squared norm of point evaluation,
\[
 \sum_j|g_{m,j}(x)|^2
   =\sup_{\|f\|_m\leq1}|f(x)|^2,
\]
and is therefore independent of the basis. By the local finite generation part of Bergman regularization, $\beta_m$ has analytic type $(1/m,\mathcal I(m\psi))$ on relatively compact open subsets of $\Omega$. This follows from the generation of the multiplier ideal by weighted $L^2$ holomorphic functions and from local $L^2$ division; see \cite[Proposition 5.7 and Section 14.A]{DemAM11}. We use this local statement, together with the point-extension estimate in the proof below, as a standard input.

\begin{proposition}\label{prop:bergman-local-estimates}
There are constants $C_1,C_2$ independent of $m$ such that, for $B(x,r)\Subset\Omega$,
\begin{equation}\label{eq:bergman-local-estimates}
 \psi(x)-\frac{C_1}{m}
 \leq\beta_m(x)
 \leq\sup_{B(x,r)}\psi+\frac1m\log\frac{C_2}{r^{2n}}.
\end{equation}
Consequently $\beta_m\to\psi$ pointwise and in $L^1_{\mathrm{loc}}$, and
\begin{equation}\label{eq:bergman-lelong-bounds}
 \max\{\nu(\psi,x)-n/m,0\}
 \leq\nu(\beta_m,x)\leq\nu(\psi,x).
\end{equation}
\end{proposition}
\begin{proof}
For $\|f\|_m\leq1$, the submean inequality gives
\[
 |f(x)|^2
 \leq C_2r^{-2n}\int_{B(x,r)}|f|^2\,\mathrm{d}\lambda
 \leq C_2r^{-2n}\mathrm{e}^{m\sup_{B(x,r)}\psi}.
\]
Taking the supremum over the unit ball of $\mathcal H_m$ and then the logarithm proves the upper bound.

For the lower bound, suppose that $\psi(x)>-\infty$. The local Ohsawa--Takegoshi point-extension theorem gives a holomorphic $f$ with $f(x)=1$ and
\[
 \int_\Omega|f|^2\mathrm{e}^{-m\psi}\,\mathrm{d}\lambda
 \leq C \mathrm{e}^{-m\psi(x)},
\]
where the constant is independent of the plurisubharmonic weight and of $m$. After normalizing $f$ in $\mathcal H_m$, we see that the squared norm of point evaluation is at least $C^{-1}\mathrm{e}^{m\psi(x)}$, which proves the lower bound. At a pole of $\psi$ the lower bound holds trivially.

Choosing $r=1/m$ in the upper bound and using upper semicontinuity of $\psi$ proves pointwise convergence. On a fixed relatively compact set, the same estimates bound $\beta_m$ above by a constant and below by $\psi-C_1/m$. Since $\psi$ is locally integrable, dominated convergence also gives convergence in $L^1_{\mathrm{loc}}$.

For the Lelong estimate, take suprema over $B(x,r)$ in the two inequalities, bounding the right-hand side of the upper one by a supremum over $B(x,2r)$. Divide by $\log r^2<0$ and let $r\downarrow0$. After the division, the term $-(2n/m)\log r$ contributes $-n/m$, while the constant terms disappear. This proves the two inequalities in \eqref{eq:bergman-lelong-bounds}; since Lelong numbers are nonnegative, the lower bound may be replaced by its maximum with zero.
\end{proof}

In \cite{DemAM11}, the local weight is $\mathrm{e}^{-2m\phi}$ and Lelong numbers are normalized by $\log|z|$. With $\psi=2\phi$, these conventions give exactly the formulas of \cref{prop:bergman-local-estimates}; in particular, the loss is $n/m$ for the weight $\mathrm{e}^{-m\psi}$ and the normalization by $\log|z|^2$.

The same construction and estimates hold for every real $m>0$. The coefficient $1/m$ of the resulting logarithmic presentation is rational exactly when $m$ is rational, and we speak of an analytic type only in that case. The global sequence constructed below is indexed by integers.

For a quasi-plurisubharmonic function $u$, the generators of $\mathcal I(mu)$ also control the multiplier ideals of $u$ at nearby weights; a small loss in the weight compensates for the finite Bergman index.

\begin{lemma}\label{lem:bergman-controlled-ideals}\index{Multiplier ideal!controlled approximation}
Let $u$ be quasi-plurisubharmonic, let $m>0$, and let $a_m$ be a function which locally satisfies $a_m=m^{-1}\log\sum_k|g_k|^2+O(1)$, where the $g_k$ generate $\mathcal I(mu)$. For $t>0$ and $s>1$,
\[
 \mathcal I(sta_m)\subseteq\mathcal I(tu)
 \qquad\text{whenever}\qquad m\geq\frac{st}{s-1}.
\]
Moreover, $u\preceq a_m$ locally, and hence $\mathcal I(tu)\subseteq\mathcal I(ta_m)$ for every $t>0$.
\end{lemma}
\begin{proof}
On a sufficiently small coordinate neighborhood, choose generators $g_1,\ldots,g_N$ of $\mathcal I(mu)$ and put $B=(\sum_k|g_k|^2)\mathrm{e}^{-mu}$. After shrinking the neighborhood, $B$ is integrable and $a_m=m^{-1}\log\sum_k|g_k|^2+O(1)$ with a bounded remainder. If $f\in\mathcal I(sta_m)$, shrink once more so that $f$ is bounded and $A=|f|^2\mathrm{e}^{-sta_m}$ is integrable. Then, almost everywhere,
\[
 |f|^2\mathrm{e}^{-tu}
 \leq C|f|^{2-2/s}A^{1/s}B^{t/m}.
\]
H\"older's inequality with exponents $s$ and $s/(s-1)$ applies because $st/(m(s-1))\leq1$, and it shows that the left side is locally integrable, also in the case of equality in the bound for $m$. Finally, the local Bergman lower estimate of \cref{prop:bergman-local-estimates} gives $u\preceq a_m$ locally; bounded smooth corrections affect neither this comparison nor any of the ideals.
\end{proof}

\Needspace{23\baselineskip}
\section{The global curvature bound}

We return to the compact manifold $X$ with its Hermitian form $\omega$. Demailly's global regularization theorem glues the local Bergman constructions while controlling the curvature error. We use it as a classical input in the following form.

\begin{theorem}[Demailly regularization]\label{thm:demailly-hermitian-input}\index{Regularization!analytic singularities}
Let $u$ be quasi-plurisubharmonic on $X$ and let $\gamma$ be a continuous real $(1,1)$-form with $\ddc u\geq\gamma$. There are quasi-plurisubharmonic functions $u_m$, defined for all sufficiently large positive integers $m$, and numbers $\epsilon_m\downarrow0$ such that
\begin{enumerate}
\item $u_m$ has analytic type $(1/m,\mathcal I(mu))$;
\item $u_m\geq u$ and $u_m\to u$ pointwise and in $L^1(X)$;
\item $\ddc u_m\geq\gamma-\epsilon_m\omega$;
\item $\max\{\nu(u,x)-n/m,0\}\leq\nu(u_m,x)\leq\nu(u,x)$ for all $x$.
\end{enumerate}
The sequence $(u_m)$ need not be monotone.
\end{theorem}

This is the form of \cite[Theorem 14.12]{DemAM11} used in this book. The local generators in the construction are the weighted Bergman generators of $\mathcal I(mu)$; smooth local corrections do not change this ideal. On overlaps the local Bergman sums are comparable, so the glued logarithm keeps the analytic type $(1/m,\mathcal I(mu))$. The construction freezes the continuous lower bound $\gamma$ on sufficiently small coordinate balls, and its curvature error combines the modulus of continuity of $\gamma$ with an overlap term that tends to zero as $m$ increases. It is this global curvature estimate, and not merely the local convergence of \cref{prop:bergman-local-estimates}, that we take from the theorem.

\begin{corollary}\label{cor:nonclosed-analytic-approximation}
Let $\theta$ be a smooth real $(1,1)$-form on $X$ and $u\in\PSH(X,\theta)$. The approximants $u_m$ of \cref{thm:demailly-hermitian-input} for $\gamma=-\theta$ satisfy
\[
 \theta+\ddc u_m\geq-\epsilon_m\omega.
\]
In particular, they have analytic type $(1/m,\mathcal I(mu))$.
\end{corollary}
\begin{proof}
Apply \cref{thm:demailly-hermitian-input} to $u$ with the lower bound $\gamma=-\theta$; the theorem does not require $\gamma$ to be closed. Adding $\theta$ to the curvature inequality gives the assertion. In the normalization of \cite{DemAM11}, the theorem is applied to the potential $u/2$, since $\frac{\mathrm{i}}{\pi}\partial\bar\partial(u/2)=\ddc u$.
\end{proof}

Thus the approximants keep the curvature lower bound up to an error that tends to zero. Their multiplier ideals determine the limit of their ordinary envelopes; see \cref{thm:dyadic-analytic-envelope-limit}.


\subsection{Regularization of closed currents}
\label{subsec:psh-closed-regularization}

We also use Demailly's regularization of closed currents, in a smooth and in an analytic form; the sources are given in \cref{sec:history-psh}. Let $T=\eta+\ddc v$ be a closed real $(1,1)$-current, where $\eta$ is a smooth closed form and $v$ is quasi-plurisubharmonic, and set $\nu(T,x)=\nu(v,x)$. Suppose that $T\geq\gamma$ for a smooth real $(1,1)$-form $\gamma$. Smooth regularization gives smooth representatives $T_j=\eta+\ddc v_j$ with
\[
 T_j\geq\gamma-C\lambda_j\omega,
\]
where $C>0$ is a fixed constant and the continuous functions $\lambda_j$ decrease pointwise to $\nu(T,\cdot)$. If all Lelong numbers of $T$ vanish, Dini's theorem shows that $\lambda_j\to0$ uniformly, so the loss of positivity becomes uniformly small. Analytic regularization \cite[Corollary~14.13]{DemAM11} instead gives representatives $S_j$ with analytic singularities and
\[
 S_j\geq\gamma-\varepsilon_j\omega,\qquad
 \varepsilon_j\downarrow0,\qquad
 \nu(S_j,x)\leq\nu(T,x).
\]
Thus analytic regularization preserves strict positivity up to a small loss and creates no positive Lelong number at a point where the Lelong number of $T$ vanishes.

\subsection{Smooth remainders after resolution}

For arguments on a modification, the bounded remainder in an analytic presentation is sometimes insufficient. The following version of regularization gives a smooth remainder after resolution.

\Needspace{13\baselineskip}
\begin{theorem}[Demailly regularization]\label{thm:resolvable-regularization-input}
Let $\alpha$ be a smooth real $(1,1)$-form and let $u\in\PSH(X,\alpha)$. There are numbers $\delta_j\downarrow0$ and functions $\widetilde u_j\in\PSH(X,\alpha+\delta_j\omega)$ decreasing to $u$ with the following properties. Each $\widetilde u_j$ is smooth outside an analytic pole set, and there is a modification $\pi_j\colon Y_j\to X$, obtained by a finite sequence of blowups with smooth centers, such that locally on $Y_j$
\[
 \pi_j^*\widetilde u_j=\sum_\ell a_\ell\log|s_\ell|^2+b,
 \qquad a_\ell>0,\quad b\in C^\infty.
\]
Here the $s_\ell$ are local holomorphic defining functions of divisors, the $a_\ell>0$ are real numbers, and $b$ is a smooth real function. The divisors may be taken to have simple normal crossings.
\end{theorem}

This classical theorem is applied to the closed current $\ddc u$ with the smooth lower bound $-\alpha$; see \cref{sec:history-psh} for its source. The lower bound need not be closed, and $X$ need not be K\"ahler. What matters here is the smooth remainder $b$ after resolution; an analytic type, whose remainder is merely bounded, would not suffice. The sequence $(\widetilde u_j)$ need not be the multiplier-ideal sequence of \cref{thm:demailly-hermitian-input}, and we do not claim that the conclusions of the two theorems hold for one and the same sequence. If $\theta+\ddc u\geq\varepsilon\omega$ for some $\varepsilon>0$, applying the theorem with $\alpha=\theta-\varepsilon\omega$ and discarding the terms with $\delta_j\geq\varepsilon$ gives approximants with $\theta+\ddc\widetilde u_j\geq(\varepsilon-\delta_j)\omega>0$; thus strict positivity is preserved.

\section{Local properties of multiplier ideals}
\label{sec:psh-local-ideals}

Multiplier ideals compare singularities through integrability. Their local valuative characterization and strong openness, recalled in \cite{SGPT}, make this comparison compatible with maxima, addition and increasing limits.

\begin{theorem}\label{thm:local-i-inputs}
For germs of plurisubharmonic functions, the following facts hold.
\begin{enumerate}
\item Equality of multiplier ideals at every positive weight is equivalent to equality of the Lelong numbers of the pullbacks at every point on every local modification. Inclusion of the ideals at every positive weight is equivalent to the reverse inequalities for these Lelong numbers.
\item If $p_j$ is an increasing sequence of plurisubharmonic functions converging almost everywhere to a plurisubharmonic function $p$, then, for each point $x$ and each $\lambda>0$,
\[
 \mathcal I(\lambda p_j)_x=\mathcal I(\lambda p)_x
 \quad\text{for all sufficiently large }j.
\]
\end{enumerate}
\end{theorem}

In the first item one works over a coordinate neighborhood of the germ, allowing smaller neighborhoods as necessary. No assertion that every local modification extends to a modification of the compact manifold is needed. The second item is the increasing-limit form of strong openness used in \cite{SGPT}. Smooth additions do not change multiplier ideals or Lelong numbers. Thus these inputs apply to quasi-plurisubharmonic functions after the local correction of \cref{prop:local-correction}.

\begin{lemma}\label{lem:hermitian-i-calculus}
Let $u,u',v,v'$ be quasi-plurisubharmonic functions on $X$. Suppose $u\sim_{\mathcal I}v$ and $u'\sim_{\mathcal I}v'$. Then
\[
 \max(u,u')\sim_{\mathcal I}\max(v,v'),\qquad
 u+u'\sim_{\mathcal I}v+v',\qquad
 a u\sim_{\mathcal I}a v\quad(a>0).
\]
These statements concern quasi-plurisubharmonic functions and require no background to be closed.
\end{lemma}
\begin{proof}
Fix a coordinate neighborhood $U\subset X$ and a smooth local modification $\pi:Y\to U$. Write $\nu(w,y)$ for the Lelong number of a quasi-plurisubharmonic function $w$ at $y$. At every $y\in Y$, the elementary Lelong identities give
\[
 \begin{aligned}
 \nu(\max(u,u')\circ\pi,y)
   &=\min\{\nu(u\circ\pi,y),\nu(u'\circ\pi,y)\},\\
 \nu((u+u')\circ\pi,y)
   &=\nu(u\circ\pi,y)+\nu(u'\circ\pi,y),\\
 \nu(a u\circ\pi,y)&=a\nu(u\circ\pi,y).
 \end{aligned}
\]
The corresponding expressions for $v,v'$ are equal by the first local input. Apply the converse direction of that input on each neighborhood. The Lelong identities are the usual local maximum, addition and scaling identities, unchanged by smooth corrections.
\end{proof}


Let $v,w$ be local quasi-plurisubharmonic functions and let $Y$ be a complex submanifold on which $v|_Y$ is genuine. Restriction and subadditivity give the two inclusions recalled in \cite{SGPT} and \cite[Section 15]{DemAM11}:
\begin{equation}\label{eq:trace-local-ideal-inputs}
 \mathcal I(v|_Y)\subseteq\mathcal I(v)\mathcal O_Y,
 \qquad
 \mathcal I(v+w)\subseteq\mathcal I(v)\mathcal I(w).
\end{equation}
The first follows from local Ohsawa--Takegoshi extension. For the second, weighted holomorphic functions on a product form the Hilbert tensor product of the two weighted spaces; its generators give the product multiplier ideal. Restriction to the diagonal then gives the inclusion. These are local results for psh functions. Smooth corrections extend them to quasi-plurisubharmonic functions with the weight convention of this book. No compact K\"ahler hypothesis is involved.


For a holomorphic submersion, multiplier ideals commute with pullback. Extension of an ideal means the ideal generated by its pullbacks.

\begin{lemma}\label{lem:submersion-multiplier-pullback}
If $X,M$ are complex manifolds, $p:X\to M$ is a holomorphic submersion and $U\in\QPSH(M)$, then for every $t>0$,
\[
 \mathcal I(tU\circ p)
       =\mathcal I(tU)\mathcal O_X.
\]
\end{lemma}
\begin{proof}
Work in product coordinates $(z,\zeta)$ with $p(z,\zeta)=z$. Write $\alpha$ for a multi-index in the fiber variables and $\mathrm{d}V_\zeta$ for their Euclidean volume. A germ in the ideal on the right is a finite sum of bounded holomorphic multiples of pullbacks of germs integrable against $\mathrm{e}^{-tU(z)}$. Fubini proves its weighted integrability on a sufficiently small product neighborhood.

Conversely, let $F(z,\zeta)$ be a weighted-integrable holomorphic germ and expand it in the fiber variables,
\[
 F(z,\zeta)=\sum_\alpha f_\alpha(z)\zeta^\alpha.
\]
The coefficient estimate for a holomorphic function on a polydisc gives, on smaller product neighborhoods,
\[
 |f_\alpha(z)|^2\leq C_\alpha
       \int |F(z,\zeta)|^2\,\mathrm{d}V_\zeta.
\]
Multiply by $\mathrm{e}^{-tU(z)}$ and integrate. Each coefficient germ belongs to $\mathcal I(tU)$. Let $J$ be the extension of this ideal to the local analytic ring $R$ of the product. Every finite Taylor sum belongs to $J$, and the remainder of total fiber degree at least $N$ belongs to $(\zeta)^N$. Thus the class of $F$ in $R/J$ lies in $(\zeta)^N(R/J)$ for every $N$. The Krull intersection theorem for the finite module $R/J$ over the Noetherian local analytic ring gives intersection zero. Hence $F\in J$. This is the usual local algebra form of convergence for coherent ideals. Smooth corrections and smooth positive volume densities do not change the integrability ideals.
\end{proof}

\section{Lelong numbers and divisorial estimates}
\label{sec:psh-divisorial-estimates}

Lelong numbers measure logarithmic concentration. Divisorial removal turns their lower bounds into inequalities of currents, while the local mass formula controls them under convergence and in a fixed background. For a prime divisor $E$ on a complex manifold, let $[E]$ be its integration current and $\nu(w,E)$ the generic Lelong number of a quasi-plurisubharmonic function $w$ along $E$.

We use the following precise local consequence of Siu's decomposition, as recalled in \cite{SGPT} and the standard current theory of \cite{DemAM11}. A closed positive $(1,1)$-current has, on a prime hypersurface $E$, coefficient-measure restriction equal to its generic Lelong number times $[E]$. In particular the remaining current after removal of any finite set of its divisorial components is positive. Distinct prime hypersurfaces do not contribute overlapping mass: their intersections have codimension at least two. These are statements about the local closed current.

\Needspace{8\baselineskip}
\begin{lemma}\label{lem:nonclosed-divisorial-removal}
Let $M$ be a complex manifold, let $\beta$ be a smooth real $(1,1)$-form on $M$, let $w\in\PSH(M,\beta)$, and let $E_1,\ldots,E_N$ be distinct prime divisors on $M$, with $N\geq1$. Then
\[
 \beta+\ddc w-\sum_{j=1}^N\nu(w,E_j)[E_j]\geq0.
\]
The form $\beta$ need not be closed.
\end{lemma}
\begin{proof}
Work on a coordinate ball and choose a smooth $\rho$ with $\ddc\rho\geq\beta$. The current
\[
 T=\ddc(w+\rho)=(\beta+\ddc w)+(\ddc\rho-\beta)
\]
is closed and positive. A smooth form, viewed through its coefficient measures, assigns zero mass to every proper analytic set. Thus $T$ and $\beta+\ddc w$ have the same restriction to $\bigcup_jE_j$. The stated local Siu fact identifies this restriction as $\sum_j\nu(w,E_j)[E_j]$. Subtracting it from the positive current $\beta+\ddc w$ leaves its coefficient-measure restriction to the complement, which is positive. This proves the inequality locally, hence globally. The argument does not identify $\beta+\ddc w$ itself as a closed current.
\end{proof}

\Needspace{12\baselineskip}
A potential whose curvature is a divisorial current plus a smooth form is determined, up to ordinary domination, by its divisorial coefficients. The comparison extends to countably many components of finite total mass.

\begin{lemma}\label{lem:divisorial-potential-domination}
Let $Y$ be a connected compact complex manifold. Let $D_j$ be distinct prime divisors, $c_j>0$ real coefficients such that $S=\sum_jc_j[D_j]$ is a positive current of finite mass, and $\beta$ a smooth real $(1,1)$-form; $[D_j]$ denotes the current of integration. Let $v\in\QPSH(Y)$ satisfy
\[
 \ddc v=S-\beta.
\]
For $a\in\QPSH(Y)$, the following conditions are equivalent:
\[
 a\preceq v,\qquad a\preceq_{\mathcal I}v,\qquad
 \nu(a,D_j)\geq c_j\quad\hbox{for every }j.
\]
\end{lemma}
\begin{proof}
The coefficient-measure restriction of $S$ to $D_j$ is $c_j[D_j]$. Distinct divisors have intersections of codimension at least two, which carry no mass of either divisor current. Countable summation preserves this assertion. Smooth forms give every divisor zero coefficient mass. Local Siu decomposition after a smooth correction therefore gives $\nu(v,D_j)=c_j$.

Ordinary domination implies ideal domination by integrability, and ideal domination implies the asserted divisorial Lelong inequalities by \cref{thm:local-i-inputs}. It remains to prove that these inequalities imply ordinary domination.

Choose a Hermitian form $\Lambda$ with $a\in\PSH(Y,\Lambda)$. By \cref{lem:nonclosed-divisorial-removal}, the divisorial inequalities allow removal of each finite partial sum of $S$ from $\Lambda+\ddc a$ while retaining positivity. Passing to the increasing finite sums gives
\[
 \Lambda+\ddc a-S\geq0.
\]
Consequently $\ddc(a-v)\geq\beta-\Lambda$. The $L^1$ difference has a quasi-plurisubharmonic representative and is bounded above on compact $Y$. This yields $a\leq v+C$ almost everywhere. \Cref{lem:qpsh-almost-everywhere-order} extends this inequality everywhere. This proves ordinary domination.
\end{proof}


\begin{lemma}\label{lem:lelong-lone-semicontinuity}
Let $u_j,u$ be quasi-plurisubharmonic on a complex manifold, with a common smooth lower curvature bound. If $u_j\to u$ in $L^1_{\mathrm{loc}}$, then at every point $x$,
\[
 \nu(u,x)\geq\limsup_j\nu(u_j,x).
\]
\end{lemma}
\begin{proof}
On a coordinate ball about $x$, add one smooth function to all the potentials to make them plurisubharmonic. Write the resulting potentials as $p_j,p$, let $n$ be the dimension of the chart, and let $\kappa$ be its Euclidean form. Distributional convergence gives weak convergence of the positive trace measures $\ddc p_j\wedge\kappa^{n-1}$ to $\ddc p\wedge\kappa^{n-1}$. For every sufficiently small radius $r$ whose boundary sphere has zero mass for the latter measure, the masses of the balls converge. The Lelong monotonicity formula bounds each $\nu(p_j,x)$ by its normalized trace mass on that ball. Taking the upper limit in $j$, and then letting such radii decrease to zero, gives $\limsup_j\nu(p_j,x)\leq\nu(p,x)$. Smooth additions do not change Lelong numbers.
\end{proof}

\begin{lemma}\label{lem:gauduchon-uniform-lelong}
Let $Z$ be connected compact complex and let $\beta$ be a smooth real $(1,1)$-form with $\PSH(Z,\beta)\ne\varnothing$. There is a finite constant $C_\beta$ such that
\[
 \nu(v,z)\leq C_\beta\qquad
 \text{for every }v\in\PSH(Z,\beta)\text{ and every }z\in Z.
\]
In particular all generic divisorial Lelong numbers in this fixed background are uniformly bounded. Additive normalization of the potentials is unnecessary.
\end{lemma}
\begin{proof}
Subtract the supremum from $v$. By \cref{thm:normalized-compactness}, these normalized potentials have a common $L^1$ bound. Cover $Z$ by finitely many nested coordinate balls. On each larger ball choose a fixed smooth quadratic function $r$ with $\ddc r\geq\beta$, so $p=v+r$ is plurisubharmonic there. The functions $p$ have uniformly bounded local $L^1$ norms.

Write $d=\dim Z$ and let $\kappa$ be the Euclidean form in the chart. Choose a nonnegative smooth cutoff $\chi$ supported in the larger ball and equal to one on an intermediate ball. Positivity and distributional integration by parts give
\[
 \int_{\text{intermediate ball}}\ddc p\wedge\kappa^{d-1}
 \leq\int\chi\,\ddc p\wedge\kappa^{d-1}
 =\int p\,\ddc\chi\wedge\kappa^{d-1}
 \leq C\|p\|_{L^1}.
\]
Every point of the smaller ball admits a surrounding Euclidean ball of a fixed radius inside the intermediate ball. The local Lelong monotonicity formula, recalled in \cite{SGPT}, bounds the Lelong number at its center by a fixed multiple of this trace mass. The finite cover gives one bound on $Z$. Smooth additions and additive normalization do not change Lelong numbers, proving the assertion for the original potentials.
\end{proof}


The Bergman estimates also control Lelong numbers on a fixed local modification. A \emph{local divisor over $X$} means a prime divisor $E$ on a smooth modification $\pi:Z\to U$ of a coordinate neighborhood. For a quasi-plurisubharmonic function $u$ on $X$, let $\nu(u\circ\pi,E)$ be its generic Lelong number along $E$, let $\operatorname{ord}_E$ denote divisorial vanishing order, and let $\operatorname{Jac}\pi$ be the local holomorphic Jacobian determinant. Write
\[
 \nu_E(u)=\nu(u\circ\pi,E),\qquad
 A_U(E)=1+\operatorname{ord}_E(\operatorname{Jac}\pi).
\]
The second number is the log discrepancy. These local modifications are not required to extend over the compact manifold.

\begin{lemma}\label{lem:bergman-divisorial-error}
For a quasi-plurisubharmonic function $u$, every fixed local divisor $E$, and every real $m>0$, one has
\[
 \max\!\left\{\nu_E(u)-\frac{A_U(E)}m,0\right\}
 \leq \frac1m\operatorname{ord}_E\mathcal I(mu)
 \leq \nu_E(u).
\]
Consequently, a local logarithmic presentation $a_m=m^{-1}\log\sum_k|f_{m,k}|^2+O(1)$, where the $f_{m,k}$ generate $\mathcal I(mu)$, satisfies $\nu_E(a_m)\to\nu_E(u)$ as $m\to\infty$.
\end{lemma}
\begin{proof}
The upper estimate follows from the local Bergman lower bound in \cref{prop:bergman-local-estimates}. The logarithm of local generators of $\mathcal I(mu)$, divided by $m$, is at least $u$ up to a locally bounded term. Pulling back and taking a generic Lelong number reverses this inequality.

For the lower estimate, take a local generator $f$ of the ideal and set $k=\operatorname{ord}_E(f)$. At a generic smooth point of $E$, choose a coordinate $z_1$ defining it. Divisorial removal, after a smooth correction, gives
\[
 u\circ\pi\leq\nu_E(u)\log|z_1|^2+C.
\]
\Needspace{9\baselineskip}
Shrink around a point where the residual factors of $f\circ\pi$ and of the Jacobian are nonvanishing. The change-of-variables formula and local integrability of $|f|^2\mathrm{e}^{-mu}$ then require integrability in $z_1$ of
\[
 |z_1|^{2(k+A_U(E)-1-m\nu_E(u))}.
\]
This implies $k+A_U(E)>m\nu_E(u)$. In particular the weaker non-strict lower bound in the statement holds. Take the minimum over a finite set of local generators and use nonnegativity of the order. Since $\nu_E(a_m)=m^{-1}\operatorname{ord}_E\mathcal I(mu)$, the last assertion follows by letting $m\to\infty$.
\end{proof}

For rational $m>0$, the presentation $a_m$ has analytic singularities in the convention of the source book. The real-parameter estimate does not require this restriction.


\begin{lemma}\label{lem:analytic-ideal-dominance}
Let $v$ be quasi-plurisubharmonic on a compact complex manifold, and let $a$ have analytic singularities. If $v\preceq_{\mathcal I}a$, then $v\preceq a$.
\end{lemma}
\begin{proof}
Resolve the coherent ideal of an analytic presentation of $a$ by a smooth modification $\pi:Y\to X$. In local coordinates $z_1,\ldots,z_n$ on $Y$, its expression has coefficients $b_i\geq0$, $1\leq i\leq r\leq n$, and a locally bounded remainder $B$:
\[
 a\circ\pi=\sum_{i=1}^r b_i\log|z_i|^2+B,
 \qquad b_i\geq0,\qquad B\in L^\infty_{\mathrm{loc}}.
\]
By \cref{thm:local-i-inputs}, the generic Lelong number of $v\circ\pi$ along $\{z_i=0\}$ is at least $b_i$. Add a smooth local function to $v\circ\pi$ to make it plurisubharmonic. Divisorial removal for this closed positive complex Hessian shows that
\[
 v\circ\pi-\sum_i b_i\log|z_i|^2
\]
has a quasi-plurisubharmonic representative. Here the difference is first interpreted off the divisor, or as an integrable distribution. Such a representative is locally bounded above. Since $B$ is locally bounded, $v\circ\pi\leq a\circ\pi+C$ on each relatively compact resolution chart away from the divisor. A finite cover above a smaller coordinate neighborhood makes $C$ uniform there. Push the inequality down off the analytic exceptional set; the canonical quasi-plurisubharmonic representatives extend it everywhere. Finitely many neighborhoods on the compact manifold give $v\leq a+C$ globally.

The only resolution input is the standard principalization of an analytic ideal. Divisorial removal is applied to a locally closed positive Hessian after smooth correction, not to a purported closed background on the original manifold.
\end{proof}


\Needspace{12\baselineskip}
\begin{lemma}\label{lem:analytic-minorant-product-ideals}
Let $u,v$ be quasi-plurisubharmonic on a compact complex manifold, and let $s$ have analytic singularities. If $u+v\leq s+O(1)$, then for analytic potentials $U_m,V_m$ of types $(1/m,\mathcal I(mu))$, $(1/m,\mathcal I(mv))$, there are numbers $\delta_m\downarrow0$ such that
\[
 U_m+V_m\leq(1-\delta_m)s+O(1)
 \quad\hbox{for all sufficiently large integers }m.
\]
The bounded error may depend on $m$.
\end{lemma}
\begin{proof}
Resolve the ideal of $s$ by a smooth modification $\pi:\widetilde X\to X$. Let $E_i$ be the components of its resolved divisor with positive coefficients $c_i$. In resolution coordinates $z_i$ defining these components, its singular part is $\sum_i c_i\log|z_i|^2$. Set
$A_i=1+\operatorname{ord}_{E_i}(\operatorname{Jac}\pi)$.
The hypothesis and \cref{lem:bergman-divisorial-error} give
\[
 \nu(U_m\circ\pi,E_i)+\nu(V_m\circ\pi,E_i)
 \geq\nu(u\circ\pi,E_i)+\nu(v\circ\pi,E_i)-\frac{2A_i}{m}
 \geq c_i-\frac{2A_i}{m}.
\]
Choose a positive integer $C\geq\max_i2A_i/c_i$ and put $\delta_m=C/m<1$ for large $m$. After a smooth local correction, divisorial removal from the positive Hessian of $(U_m+V_m)\circ\pi$ shows that
\[
 (U_m+V_m)\circ\pi-(1-\delta_m)\sum_i c_i\log|z_i|^2
\]
has a quasi-plurisubharmonic representative. It is locally bounded above. A finite resolution cover and proper descent therefore give
$U_m+V_m\leq(1-\delta_m)s+O(1)$, just as in \cref{lem:analytic-ideal-dominance}. The bounded constant may depend on $m$. If the divisor is empty, $s$ is bounded and upper bounds for $U_m,V_m$ prove the assertion directly.
\end{proof}


\Needspace{14\baselineskip}
\section{Potential regularity and metric distance}
\label{sec:psh-holder-geometry}

A H\"older modulus of continuity of the potential controls the trace mass of a positive current on small balls. Averaging the lengths of segments against this trace measure then bounds the distance of the metric defined by the current on its regular locus. This turns regularity of the potential into a metric estimate without any pointwise bound on the Hessian.

Let $U\subset\C^n$ be a coordinate domain with Euclidean form $\omega_{\mathrm E}=\ddc|z|^2$ and Euclidean volume $\mathrm{d}\lambda$. Let $0<\alpha\leq1$, let $\theta$ be a smooth real $(1,1)$-form on $U$, and let $v\in C^{0,\alpha}(U)$ satisfy $T=\theta+\ddc v\geq0$. Write $\mu=T\wedge\omega_{\mathrm E}^{n-1}$ for the trace measure of $T$.

\begin{theorem}\label{thm:holder-current-distance}\index{Holder potential@H\"older potential!trace mass}\index{Holder potential@H\"older potential!metric distance}
For every open set $V\Subset U$ there are constants $C,r_0>0$ such that
\begin{equation}\label{eq:trace-ball}
 \mu(B(x,r))\leq C r^{2n-2+\alpha},\qquad x\in V,\quad0<r<r_0.
\end{equation}
The total variation of every coefficient measure of $T$ is bounded by a constant multiple of $\mu$. Hence the coefficient measures of $T$ vanish on sets of zero Hausdorff $(2n-2+\alpha)$-measure, and all Lelong numbers of $v$ vanish. The estimate implies neither absolute continuity of $\mu$ nor a pointwise bound for the Hessian of $v$.

Suppose in addition that $T$ is a smooth Hermitian form $\eta$ on $U\setminus A$ for a proper analytic subset $A\subset U$. For $x,y\in U\setminus A$, let $d_\eta^U(x,y)$ be the infimum of the $\eta$-lengths of piecewise smooth paths in $U\setminus A$ joining $x$ and $y$. Then every point of $U$ has a neighborhood $V\Subset U$ such that, for a constant $C$,
\begin{equation}\label{eq:holder-distance}
 d_\eta^U(x,y)\leq C|x-y|^{\alpha/2},\qquad x,y\in V\setminus A.
\end{equation}
In particular, all sequences in $V\setminus A$ converging to a given point of $A\cap V$ are Cauchy for $d_\eta^U$ and have a common limit in the metric completion. If these local hypotheses hold in charts covering $X$, with a proper analytic subset $A\subset X$ and a Hermitian metric $\eta$ on $X\setminus A$, then the metric completion of $X\setminus A$ for its intrinsic $\eta$-distance is compact, and the identity of $X\setminus A$ extends to a continuous surjection of $X$ onto it. The metric $\eta$ is incomplete when $A\ne\varnothing$. Both exponents in \eqref{eq:trace-ball} and \eqref{eq:holder-distance} are sharp.
\end{theorem}

\begin{proof}
Choose a nonnegative cutoff function $\chi$ equal to one on $B(x,r)$ and supported in $B(x,2r)$, with second derivatives bounded by $Cr^{-2}$. Positivity and integration by parts, after subtracting the constant $v(x)$ from $v$, give
\[
 \mu(B(x,r))\leq C_1r^{2n}+C_2r^{2n-2}
  \operatorname{osc}_{B(x,2r)}v
 \leq Cr^{2n-2+\alpha}.
\]
Here we used that the constant $v(x)$ integrates to zero against $\ddc\chi\wedge\omega_{\mathrm E}^{n-1}$. The matrix $(a_{ij})$ of Radon--Nikodym derivatives of the coefficient measures with respect to $\mu$ is positive semidefinite almost everywhere, so $|a_{ij}|\leq(a_{ii}+a_{jj})/2$. A ball cover proves the assertion about Hausdorff null sets. After a smooth local correction making $v$ plurisubharmonic, dividing \eqref{eq:trace-ball} by $r^{2n-2}$ shows that all Lelong numbers vanish.

For the distance estimate, let $m=2n$ and let $\tau$ be the Euclidean trace of $\eta$. Then $\tau\,\mathrm{d}\lambda\leq C\mu$ on the regular locus and $|\xi|_\eta\leq C\sqrt\tau|\xi|_{\mathrm E}$ for tangent vectors $\xi$. Polar integration bounds the average $\eta$-length of the segments from $x$ to the points of $B(x,R)$ by a constant times $\int_{B(x,R)}\sqrt\tau(\zeta)|\zeta-x|^{1-m}\,\mathrm{d}\lambda(\zeta)$. On a dyadic annulus of radius $s$, the Cauchy--Schwarz inequality and \eqref{eq:trace-ball} bound the contribution to this integral by $Cs^{1-m}s^{m/2}s^{(m-2+\alpha)/2}=Cs^{\alpha/2}$. Summing over the annuli gives the bound $CR^{\alpha/2}$.

Now let $x,y\in V\setminus A$, and average the lengths of the broken segments $[x,z]\cup[z,y]$ over $z$ in the ball of radius $|x-y|/2$ about the midpoint of $x$ and $y$. By the preceding bound, some such path has length at most $C|x-y|^{\alpha/2}$. For almost every $z$, both segments avoid $A$: proper analytic sets are locally countable unions of smooth strata of real codimension at least two, and the cones from a fixed regular point over these strata have zero volume. If $V$ is small enough, all these paths stay inside $U$.

By this estimate, all regular sequences converging to a point of $A$ have a common limit in the completion, and the identity of the regular locus extends continuously across $A$. By compactness, the image of this extension is the entire completion; different points of $A$ may be identified. A sequence converging to a point of $A$ cannot converge in the metric to a regular point, since $\eta$ is bounded below by a positive multiple of the Euclidean metric on a small closed neighborhood of that point. Hence the metric $\eta$ is incomplete when $A$ is nonempty. Finally, the local model $v=|z|^\alpha+|w|^2$, in coordinates $(z,w)\in\C\times\C^{n-1}$, has trace mass of exact order $r^{2n-2+\alpha}$ on balls centered on $\{z=0\}$, and the distance to $\{z=0\}$ is of exact order $|z|^{\alpha/2}$; this proves both sharpness claims. In later applications to potentials with logarithmic poles, the theorem is applied to the positive residual current that remains after the divisorial part has been removed.
\end{proof}

\subsection{Dini moduli and limits of distances}
\label{sec:dini-metric-distance}

The annular argument of \cref{thm:holder-current-distance} also applies to a continuous potential whose modulus of continuity is not a power; what matters is a Dini condition on the square root of the modulus. This gives uniform distance estimates for families whose H\"older constants are not controlled.

Keep the coordinate domain $U\subset\C^n$ with its Euclidean form $\omega_{\mathrm E}$ and volume $\mathrm{d}\lambda$, but now let $v\in C^0(U)$ and let $\theta$ be a smooth real $(1,1)$-form with $T=\theta+\ddc v\geq0$. On the relatively compact charts under consideration, choose a nondecreasing function $m\colon(0,r_0)\to[0,+\infty)$ with $m(r)\to0$ as $r\to0$ such that
\[
 |v(x)-v(y)|\leq m(|x-y|),\qquad |x-y|<r_0.
\]
Put
\begin{equation}\label{eq:dini-distance-modulus}
 D(r)=r+\int_0^r\frac{\sqrt{m(t)}}{t}\,\mathrm{d}t.
\end{equation}

\Needspace{14\baselineskip}
\begin{theorem}\label{thm:dini-current-distance}\index{Dini modulus!metric distance}\index{Metric completion!Dini potential}
Assume that $D(r_0/2)<+\infty$ and that $T$ is a smooth Hermitian form $\eta$ on $U\setminus A$ for a proper analytic subset $A\subset U$. Then every point of $U$ has a neighborhood $V\Subset U$ with constants $C,r_1>0$ such that, for $x,y\in V$,
\[
 d_\eta^U(x,y)\leq C D(8|x-y|),
 \qquad x,y\notin A,\quad |x-y|<r_1.
\]
The constants depend only on the charts and an upper bound for $|\theta|$; the dependence on $v$ enters only through $D$. All sequences in $V\setminus A$ converging to a given point of $A\cap V$ have a common limit in the metric completion. If these hypotheses hold on a finite cover of a compact connected manifold by charts, the metric completion of the regular locus is compact, and the identity of the regular locus extends to a continuous surjection of the manifold onto it. If $A$ is nonempty, the metric $\eta$ on the regular locus is incomplete.
\end{theorem}
\begin{proof}
Write $d=2n$ and $\mu=T\wedge\omega_{\mathrm E}^{n-1}$. The same cutoff calculation as in \cref{thm:holder-current-distance}, with a constant subtracted from $v$, gives
\[
 \mu(B(x,s))\leq C\bigl(s^d+s^{d-2}m(2s)\bigr).
\]
This uses only the bound for $\theta$; closedness of $\theta$ is not needed. If $\tau$ is the Euclidean trace of $\eta$, then $\tau\,\mathrm{d}\lambda\leq C\mu$ on the regular locus. Polar integration bounds the average length of the segments from a regular point $x$ to the points of $B(x,R)$ by
\[
 C\int_{B(x,R)}\sqrt{\tau(\zeta)}
                      |\zeta-x|^{1-d}\,\mathrm{d}\lambda(\zeta).
\]
On an annulus of outer radius $s$, the Cauchy--Schwarz inequality and the trace bound make the contribution to this integral at most
\[
 Cs^{1-d}s^{d/2}
       \bigl(s^d+s^{d-2}m(2s)\bigr)^{1/2}
 \leq C\bigl(s+\sqrt{m(2s)}\bigr).
\]
For $s_j=2^{-j}R$, monotonicity gives
\[
 \sum_{j\geq0}\sqrt{m(2s_j)}
 \leq\frac1{\log2}\int_0^{4R}\frac{\sqrt{m(t)}}t\,\mathrm{d}t.
\]
Indeed, $\sqrt{m(2s_j)}\log2$ is at most the integral of $\sqrt{m(t)}/t$ over $[2s_j,4s_j]$, and these intervals have disjoint interiors; no doubling condition on $m$ is needed. The average length is therefore at most $CD(4R)$.

For $x,y\in V\setminus A$, average the lengths of the broken segments $[x,z]\cup[z,y]$ over $z$ in the ball of radius $|x-y|/2$ about the midpoint of $x$ and $y$, applying the preceding bound with $R=2|x-y|$ at each endpoint. Since the volume of this ball is comparable to $R^d$, some broken path has length at most $CD(8|x-y|)$. As in the proof of \cref{thm:holder-current-distance}, almost every such path avoids $A$, and for small $V$ it stays inside $U$.

Since $D(r_0/2)<+\infty$, we have $D(r)\to0$ as $r\to0$. Hence any two regular sequences converging to the same point are Cauchy and have the same limit in the completion. The argument for the completion in the proof of \cref{thm:holder-current-distance} applies with this modulus. A positive lower bound for $\eta$ near a regular point prevents a sequence converging to a point of $A$ from converging in the metric to the regular point. This also proves incompleteness when $A$ is nonempty.
\end{proof}

\Needspace{10\baselineskip}
\begin{corollary}\label{cor:dini-distance-rates}\index{Dini modulus!logarithmic distance rate}
In the setting of \cref{thm:dini-current-distance}, let $V$ be as in the theorem. For $0<\beta\leq1$, a modulus $m(r)\leq Cr^\beta$ gives $d_\eta^U(x,y)\leq C'|x-y|^{\beta/2}$ for sufficiently close $x,y\in V\setminus A$. If instead
\[
 m(r)\leq C(-\log r)^{-a},\qquad a>2,
\]
then for sufficiently close $x,y\in V\setminus A$,
\[
 d_\eta^U(x,y)\leq C'(-\log|x-y|)^{1-a/2}.
\]
These are only sufficient conditions; we do not claim that the Dini condition is necessary.
\end{corollary}
\begin{proof}
Integrating the square root of the power modulus gives a multiple of $r^{\beta/2}$. For the logarithmic modulus, the substitution $s=-\log t$ gives
\[
 \int_0^r\frac{(-\log t)^{-a/2}}t\,\mathrm{d}t
 =\frac{(-\log r)^{1-a/2}}{a/2-1}.
\]
In either case the term $r$ in $D(r)$ is bounded by the asserted rate for small $r$. The theorem now gives both estimates; the fixed dilation of the argument of $D$ changes only the constants.
\end{proof}

\Needspace{14\baselineskip}
\begin{corollary}\label{cor:dini-family-distance-limits}\index{Metric distance!uniform pseudometric limit}
Let $\omega_0$ be a Hermitian form on $X$ and $d_{\omega_0}$ the distance of the associated Riemannian metric. For each $j$, let $T_j=\theta_j+\ddc v_j\geq0$, where $\theta_j$ is a smooth real $(1,1)$-form, $v_j$ is continuous, and $T_j$ is a smooth Hermitian form off a proper analytic subset $A_j$. Assume
\[
 |\theta_j|_{\omega_0}\leq\Lambda,\qquad
 |v_j(x)-v_j(y)|\leq m(d_{\omega_0}(x,y)),
\]
with one fixed $\Lambda$ and one nondecreasing modulus $m$ tending to zero such that the function in \eqref{eq:dini-distance-modulus} satisfies $D(r_0/2)<+\infty$. By \cref{thm:dini-current-distance}, the identity of $X\setminus A_j$ extends to a continuous surjection of $X$ onto the metric completion of $X\setminus A_j$; let $d_j$ be the pseudometric on $X$ obtained by pulling back the completed distance. Then the diameters of the $d_j$ are uniformly bounded, and a subsequence converges uniformly on $X\times X$ to a continuous pseudometric. If $T_j\geq c\omega_0$ on the regular loci for one $c>0$, the limit is a metric and dominates $\sqrt c\,d_{\omega_0}$.
\end{corollary}
\begin{proof}
A finite fixed chart cover and the theorem give common constants $C,r_1>0$ such that
\[
 d_j(x,y)\leq CD(Cd_{\omega_0}(x,y)),
 \qquad d_{\omega_0}(x,y)<r_1.
\]
For a pair meeting $A_j$, approximate it by regular pairs in the manifold topology. The completed distance and $D$ are continuous, so the same estimate holds for all such pairs in $X\times X$. No uniform bound on the geometry of the analytic sets enters the constants. Dividing a minimizing geodesic of $d_{\omega_0}$ into a uniformly bounded number of short pieces and using the triangle inequality gives a uniform diameter bound. The triangle inequality also gives
\[
 |d_j(x,y)-d_j(x',y')|
 \leq d_j(x,x')+d_j(y,y'),
\]
so the distance functions are equicontinuous. The Arzel\`a--Ascoli theorem gives a subsequence $d_{j_k}$ converging uniformly on $X\times X$ to a function $d_\infty$. Nonnegativity, symmetry, vanishing on the diagonal and the triangle inequality pass to the limit. The lower bound $T_j\geq c\omega_0$ gives $d_j\geq\sqrt c\,d_{\omega_0}$, first for regular pairs and then for all pairs, which proves the last assertion.

Without this lower bound the limit may degenerate. For $\varepsilon_j\downarrow0$, take $\theta_j=\varepsilon_j\omega_0$ and $v_j=0$, with $A_j=\varnothing$. These data satisfy the hypotheses with $m(r)=r^2$, and $d_j=\sqrt{\varepsilon_j}\,d_{\omega_0}$ converges uniformly to zero.
\end{proof}

\section{Obstacle envelopes}
\label{sec:obstacle-envelopes}

An obstacle selects the largest potential lying below it. The local upper-envelope theorem gives the construction once the competitors have a common upper bound; normalized compactness supplies that bound for obstacles finite almost everywhere.

A smooth real $(1,1)$-form $\theta$ is called \emph{big} if there are $\rho\in\PSH(X,\theta)$ and $\delta>0$ such that
\[
 \theta+\ddc\rho\geq\delta\omega.
\]
Here $\omega$ is any Hermitian metric. Changing that metric changes $\delta$, not the condition. Neither this definition nor the existence of $\theta$-plurisubharmonic functions requires $\theta$ to be closed.

For any smooth real $(1,1)$-form $\theta$ and an extended-real obstacle $f$ on $X$, defined at least outside a pluripolar set, write
\[
 P_\theta(f)=\left(\sup\{v\in\PSH(X,\theta):v\leq f
                      \ \hbox{q.e.}\}\right)^*.
\]
For $u,v\in\PSH(X,\theta)$, set $P_\theta(u,v)=P_\theta(\min(u,v))$. The star denotes upper semicontinuous regularization, and the empty supremum is $-\infty$. Quasi-everywhere means outside a pluripolar set. This convention also gives meaning to differences of quasi-plurisubharmonic functions at common poles. For an upper semicontinuous obstacle, the pointwise definition gives the same competitors: a quasi-plurisubharmonic function lying below it almost everywhere lies below it everywhere, by the local submean inequality and upper semicontinuity of the obstacle.

The local envelope theorem from \cite{SGPT} states that the supremum of a nonempty, locally uniformly upper-bounded family of plurisubharmonic functions has a plurisubharmonic upper semicontinuous regularization. This regularization differs from the supremum only on a pluripolar set. To retain the lower bound $-\theta$, fix a constant positive form $\kappa$ in local coordinates. For each $\epsilon>0$, shrink the neighborhood and choose a constant form $\theta'$ with $-\epsilon\kappa\leq\theta-\theta'\leq\epsilon\kappa$. A smooth potential $r$ for $\theta'+\epsilon\kappa$ makes every competitor plus $r$ plurisubharmonic. The local envelope theorem then gives $\theta+\ddc P_\theta(f)\geq-2\epsilon\kappa$ there. Letting $\epsilon$ tend to zero proves the exact bound. The same local theorem shows that regularization changes the supremum only on a pluripolar set.

Here is a useful way to check the upper bound for an unbounded obstacle. If $f\leq A$ for a real constant $A$ on a Borel set $E$ of positive $\mathrm{d}V$-measure and $h\leq f$ is a competitor, normalize $\widetilde h=h-\sup_Xh$. By \cref{thm:normalized-compactness},
\[
 (\sup_Xh-A)\operatorname{vol}_{\mathrm{d}V}(E)
 \leq\int_E|\widetilde h|\,\mathrm{d}V\leq C_\theta
 \quad\hbox{whenever }\sup_Xh>A.
\]
Thus the competitors have a common upper bound. In particular this applies to any measurable obstacle finite almost everywhere, by taking one of its finite sublevel sets of positive volume.

If $\PSH(X,\theta)$ is nonempty, its normalized envelope
\[
 V_\theta=P_\theta(0)
\]
is $\theta$-plurisubharmonic. Indeed, a normalized $\theta$-plurisubharmonic function is a competitor, and the family is bounded above by zero. Moreover,
\[
 u-\sup_Xu\leq V_\theta\qquad(u\in\PSH(X,\theta)).
\]
Thus $V_\theta$ has minimal singularities. If $\theta\geq0$, then $V_\theta=0$.

\begin{lemma}\label{lem:bounded-below-envelope}
Let $\eta$ be a smooth semipositive real $(1,1)$-form on $X$. If an obstacle $f$ is finite almost everywhere and bounded below quasi-everywhere by a constant, then $P_\eta(f)$ is a bounded $\eta$-plurisubharmonic function.
\end{lemma}
\begin{proof}
The lower constant is a competitor. The preceding finite-sublevel argument gives a common upper bound for all competitors. The local envelope theorem therefore gives a bounded envelope.
\end{proof}

\begin{proposition}
\label{prop:continuous-obstacle-envelope}
If $f$ is continuous and $\omega$ is Hermitian, then $P_\omega(f)$ is continuous. Moreover,
\[
 \|P_\omega(f)-P_\omega(g)\|_\infty\leq\|f-g\|_\infty
 \quad(f,g\in C^0(X)).
\]
\end{proposition}
\begin{proof}
The constant $\min_Xf$ is a competitor, so $p=P_\omega(f)$ is bounded. The local envelope theorem gives $p\leq f$ quasi-everywhere, and continuity of $f$ makes this inequality pointwise. Choose smooth $v_j\downarrow p$ by \cref{thm:bounded-exact-regularization}. For every $\epsilon>0$, the continuous functions $v_j-f$ decrease to a function at most zero. Compactness shows that $v_j\leq f+\epsilon$ everywhere for all large $j$: otherwise their decreasing compact $\epsilon$-superlevel sets would have nonempty intersection. Hence $v_j-\epsilon$ is a competitor and
\[
 p\leq v_j\leq p+\epsilon.
\]
Thus $v_j\to p$ uniformly, proving continuity. The last inequality follows directly from the competitors and
$P_\omega(g+c)=P_\omega(g)+c$ for constants $c$.
\end{proof}

\begin{definition}\label{def:singularity-envelope}
For $u\in\PSH(X,\theta)$, set
\[
 P_\theta[u]
 =\left(\lim_{C\to+\infty}P_\theta(u+C,0)\right)^*
 =\left(\sup\{h\in\PSH(X,\theta):h\preceq u,\ h\leq0\}\right)^*.
\]
We call $u$ an \emph{envelope fixed point} if $P_\theta[u]=u$.
\end{definition}

The two formulas agree by taking the union of the competitor families. They also agree with the definition using $V_\theta$ in place of zero: every $\theta$-plurisubharmonic competitor below zero is below $V_\theta$. The envelope is a genuine potential, since the competitors contain $u-\sup_Xu$ and are bounded above by zero. In particular, $u\preceq P_\theta[u]$. The last relation is not an assertion that the two functions differ by a bounded amount.

Directly from the competitors, $P_\theta[u]$ is unchanged when $u$ is replaced by a potential of the same ordinary singularity type, and
\[
 u\preceq v\quad\Longrightarrow\quad P_\theta[u]\leq P_\theta[v].
\]
These statements concern ordinary singularities. Replacing $\preceq$ by the multiplier-ideal order would define a different operator, which will require a separate comparison.

\Needspace{9\baselineskip}
\section{The minimum principle}
\label{sec:psh-minimum}

An infimum of plurisubharmonic functions need not be plurisubharmonic. Kiselman's principle gives a different operation: take an infimum in a real parameter whose complexification is plurisubharmonic and invariant in the imaginary direction. Freezing the background coefficients preserves the precise lower curvature bound.

\begin{lemma}\label{lem:invariant-strip-convexity}
Let $I\subset\mathbb R$ be a nonempty open interval, set $S_I=\{z\in\mathbb C:\Real z\in I\}$, and let $\pi:X\times S_I\to X$ be the projection. Suppose $U(x,z)=u_{\Real z}(x)$ is quasi-plurisubharmonic and $\pi^*\theta+\ddc U\geq0$. For each $x$, the function $t\mapsto u_t(x)$ is finite and convex or identically $-\infty$. The finite alternative holds for almost every $x$.
\end{lemma}
\begin{proof}
A local smooth correction making $U$ plurisubharmonic can be chosen to depend only on $x$. Restriction to a vertical complex line is therefore subharmonic or identically $-\infty$. In the former case its distributional Laplacian is the second real derivative of $t\mapsto u_t(x)$, which is nonnegative. The resulting convex representative is finite on $I$ and agrees everywhere with the subharmonic function by the submean inequality. Finally, $U$ is locally integrable on the product. Fubini's theorem, applied in coordinate balls times a fixed rectangle in $S_I$, excludes the identically $-\infty$ alternative for almost every $x$.
\end{proof}

\begin{proposition}\label{prop:hermitian-kiselman}
Under the hypotheses of \cref{lem:invariant-strip-convexity}, the function
\[
 q(x)=\inf_{t\in I}u_t(x)
\]
is either a genuine $\theta$-plurisubharmonic function or identically $-\infty$. No regularization of the infimum is necessary. If $I$ is bounded, the infimum is genuine.
\end{proposition}
\begin{proof}
The classical local Kiselman principle says that the infimum in the real variable of a plurisubharmonic function on a product with a strip, invariant in its imaginary variable, is plurisubharmonic or identically $-\infty$. Its parameter interval may be unbounded. We use this local result from \cite{SGPT}.

Fix $x_0\in X$, choose local holomorphic coordinates centered at $x_0$, and let $\kappa$ be their Euclidean form. Fix $\eta>0$. On a sufficiently small coordinate ball the constant form $\beta=\theta(x_0)$ satisfies
\[
 \beta-\eta\kappa\leq\theta\leq\beta+\eta\kappa.
\]
Choose a quadratic function $r$ with $\ddc r=\beta+\eta\kappa$. Then $U+r\circ\pi$ is psh on the ball times $S_I$. Its infimum is $q+r$, so the local principle gives
\[
 \ddc q\geq-\beta-\eta\kappa\geq-\theta-2\eta\kappa
\]
where the infimum is genuine. In a fixed atlas the Euclidean forms are bounded by a fixed Hermitian metric. Letting $\eta\downarrow0$ proves the exact lower bound $\ddc q\geq-\theta$. The local alternative propagates across overlapping balls on connected $X$: a genuine psh function cannot equal $-\infty$ on an open set. Thus either $q$ is genuine everywhere as a quasi-plurisubharmonic function or $q\equiv-\infty$. The infimum is upper semicontinuous, as an infimum of upper semicontinuous functions.

Suppose $I$ is bounded. By \cref{lem:invariant-strip-convexity}, the function $t\mapsto u_t(x)$ is finite and convex for almost every $x$. A supporting affine function at any interior point bounds it from below on $I$. Hence $q(x)>-\infty$ almost everywhere, excluding the identically $-\infty$ alternative.
\end{proof}

For example, the proposition applies to $u_t-t\tau$ for a fixed real $\tau$. On an unbounded interval the infimum may be $-\infty$; this alternative will encode the cutoff of a test curve.

\section{Chern-connection regularization}
\label{sec:psh-chern-regularization}\index{Regularization!Chern connection}

Fix $C_0>0$ and $u\in\PSH(X,\omega)$ with $-C_0\leq u\leq0$. Distances, gradients and tensor norms use the Riemannian metric associated to $\omega$.

Let $R_tu$ be Demailly's regularization formed with the Chern connection of $\omega$. It averages $u$ over the tangent fiber with a fixed smooth nonnegative radial kernel of integral one. The averaging map is a smooth realization of the holomorphic part of the formal fiberwise expansion of the exponential map; it need not be genuinely holomorphic in the fiber. The following classical estimate controls the truncated transform, including its mixed base--fiber terms; its sources are recorded in \cref{sec:history-psh}.

\Needspace{23\baselineskip}
\begin{lemma}\label{lem:chern-kiselman-input}
There are constants $A,K>0$ and $t_*>0$, depending only on the fixed geometry and bound for $u$, such that $R_tu$ is smooth for $0<t<t_*$,
\[
 -C_0\leq R_tu\leq0,\qquad
 R_tu+Kt^2\downarrow u\quad\text{as }t\downarrow0.
\]
For $0<\delta<t_*$ and $c>0$, the bounded function
\[
 \Phi_{c,\delta}
   =\inf_{0<t\leq\delta}
       \left\{R_tu+K(t^2-\delta^2)+K(t-\delta)
                         -c\log(t/\delta)\right\}
\]
is quasi-plurisubharmonic and satisfies
\[
 \omega+\ddc\Phi_{c,\delta}\geq-(Ac+2K\delta)\omega.
\]
The infimum is attained at some $t>0$ at every point.
\end{lemma}

Only this classical curvature estimate is retained as an input. Its linear term $K(t-\delta)$ matters in the Hermitian case: it absorbs the mixed base--fiber terms in the complex Hessian of the regularization. The function $-c\log(t/\delta)$ tends to $+\infty$ as $t\downarrow0$, while $u$ is bounded, which explains attainment away from zero. We require neither smoothness of the minimizing radius nor smoothness of the transform across the region where that radius equals $\delta$.

We record two elementary estimates for the averaging operators so that their use does not require a K\"ahler volume identity.

\Needspace{11\baselineskip}
\begin{lemma}\label{lem:chern-smoothing-integral}
For all sufficiently small $t>0$,
\[
 \|R_tu-u\|_{L^1(\mathrm{d}V)}\leq Ct^2,
 \qquad \|\nabla R_tu\|_\infty\leq C/t.
\]
The constants are uniform for $-C_0\leq u\leq0$.
\end{lemma}
\begin{proof}
Let $z$ be local base coordinates and $\xi\in\mathbb C^n$ the fiber variable in a smooth orthonormal tangent frame. Write $\rho(|\xi|^2)$ for the fixed radial averaging kernel, normalized against Lebesgue measure $\mathrm{d}\xi$. In these coordinates the averaging map $F$ gives
\[
 R_tu(z)=\int u(F(z,t\xi))\rho(|\xi|^2)\,\mathrm{d}\xi,
 \qquad F(z,0)=z,
\]
where the kernel is supported in $|\xi|\leq1$. The map $F$ is smooth in these variables and its fiber differential at zero is an isomorphism. We first claim that the adjoint with respect to $\mathrm{d}V$ satisfies
\[
 R_t^*1=1+O(t^2)
\]
uniformly. To see this, insert a smooth partition of unity in the $z$ variable with each term compactly supported in a frame chart. For fixed $\xi$, change variables $x=F(z,t\xi)$ on a neighborhood of that support. The inverse map and its Jacobian have uniform Taylor expansions in $t$. The term of order zero in the transformed density is the chosen partition function at $x$. Its term of order one is linear in $\xi$. Integrating that term against the even radial kernel gives zero. The remaining density is bounded by $Ct^2\mathrm{d}V(x)$, since the second derivatives of all maps, partition functions and volume coefficients are bounded on these finitely many compact charts. Summing the partition functions proves the claim.

Consequently
\[
 \left|\int_X(R_tu-u)\,\mathrm{d}V\right|
  =\left|\int_Xu(R_t^*1-1)\,\mathrm{d}V\right|\leq Ct^2.
\]
The preceding lemma gives $R_tu-u\geq-Kt^2$. For a real function $a\geq-Kt^2$ one has $|a|\leq a+2Kt^2$. Apply this inequality and integrate to obtain the first estimate.

For the derivative bound, change variables in the fiber to express $R_tu(z)=\int_Xu(x)k_t(z,x)\,\mathrm{d}V(x)$. In local coordinates the smooth kernel is supported where $|x-z|\leq Ct$ and satisfies
\[
 |\nabla_z k_t(z,x)|\leq Ct^{-2n-1}.
\]
Indeed, the inverse fiber map and its Jacobian have uniformly bounded derivatives; differentiating the scaled radial argument introduces at most one factor $t^{-1}$ in addition to $t^{-2n}$. Smooth compact support of the kernel causes no boundary term. Integration over a ball of volume at most $Ct^{2n}$ gives $\int|\nabla_zk_t|\,\mathrm{d}V\leq C/t$. Multiplication by the uniform bound for $u$ proves the assertion.
\end{proof}

The first estimate uses cancellation of the odd first-order term of an averaging kernel. It does not use conservation of Hermitian Monge--Amp\`ere mass.


\subsection{The general curvature modulus}
\label{subsec:general-curvature-modulus}

Let $\theta$ now be any smooth real $(1,1)$-form and let $u\in\PSH(X,\theta)\cap L^\infty(X)$. Use the same averaging operators and metric.

\Needspace{10\baselineskip}
We first specify the classical input for this background. Apply the classical Chern-connection regularization theorem to the closed current $\ddc u$ with smooth, possibly nonclosed lower bound $-\theta$. Its tangent-bundle curvature bound is dominated by a fixed multiple of $\omega$. The averaging operators have a constant $K_0>0$ such that
\[
 R_tu+K_0t^2\downarrow u.
\]
In local base coordinates $z$ and a complex radius variable $w$, put $t=|w|$. Write $\zeta\in T_z^{1,0}X$ and $\eta\in\mathbb C$ for the base and radius components of a tangent vector, and set $\lambda(z,t)=t\partial_t(R_tu(z)+K_0t^2)$. With fixed constants $A,C>0$ and an error function $\epsilon_0(t)\geq0$ tending to zero, the theorem bounds the full base--fiber complex Hessian by the base form $-\theta-A\lambda(z,t)\omega-\epsilon_0(t)\omega$, together with mixed and vertical errors bounded by
$C(|\zeta||\eta|+|\eta|^2)$. Here $\lambda$ is nonnegative and increasing in $t$, and $\epsilon_0(t)\downarrow0$ as $t\downarrow0$. This is the classical regularization estimate for a smooth lower bound.

\Needspace{13\baselineskip}
\begin{lemma}\label{lem:general-kiselman-modulus}
For $K$ sufficiently large, there is an increasing function $\epsilon(\delta)\geq0$ tending to zero at zero such that, for every small $\delta>0$ and $c>0$,
\[
 \Phi_{c,\delta}
   =\inf_{0<t\leq\delta}
      \{R_tu+K(t^2-\delta^2)+K(t-\delta)-c\log(t/\delta)\}
\]
is bounded, quasi-plurisubharmonic, and satisfies
\[
 \theta+\ddc\Phi_{c,\delta}
                       \geq-(Ac+\epsilon(\delta))\omega.
\]
The infimum is attained away from zero. One may also use $K$ in the monotone smoothing inequality above.
\end{lemma}
\begin{proof}
Take $K\geq K_0$ and put $H(z,t)=R_tu(z)+Kt^2+Kt$. The function $H$ is increasing and strictly convex as a function of $\log t$, because $R_tu+K_0t^2$ has nonnegative increasing logarithmic derivative and the added terms have positive second logarithmic derivative. Its logarithmic derivative $\widetilde\lambda$ is at least $\lambda$.

The complex Hessian of $K|w|$ has a positive vertical term proportional to $K/|w|$. For sufficiently small $|w|=t$, it absorbs the vertical error and the mixed term, using
\[
 C|\zeta||\eta|
       \leq C' t|\zeta|^2+\frac{K}{8t}|\eta|^2.
\]
After increasing constants and replacing $\lambda$ by the larger $\widetilde\lambda$, the full base--fiber Hessian of $H$ is therefore bounded below by the base form
\[
 -\theta-(A\widetilde\lambda+\epsilon(t))\omega,
\]
where an increasing majorant of $\epsilon_0(t)+C't$ may be used.

Boundedness of $u$ makes the logarithmic penalty tend uniformly to $+\infty$ at zero. Thus the minimum occurs on a compact radial interval away from zero. Strict logarithmic convexity makes the minimizing radius unique. At an interior minimum $\widetilde\lambda=c$; at the upper cap $t=\delta$, one has $\widetilde\lambda\leq c$. No differentiability of the minimizing radius is needed.

For completeness, localize near any base point and fix $a>0$. Continuity of compact minimizers and of $\widetilde\lambda$ permits a radial interval containing all nearby minimizers on which $\widetilde\lambda\leq c+a$. In the cap case keep the original upper radius $\delta$; an open interval approaching the cap has the same infimum. The lower endpoint can be chosen below all minimizers by the uniform logarithmic penalty. Monotonicity of $\widetilde\lambda$ ensures its bound throughout the chosen interval. The minimum principle for rotation-invariant plurisubharmonic functions, applied after a local quadratic correction, now gives the lower curvature bound with $A(c+a)+\epsilon(\delta)$ and an arbitrarily small error from freezing the smooth base forms. Let these errors and $a$ tend to zero. This proves the asserted current inequality and quasi-plurisubharmonicity. Finally, $K\geq K_0$ preserves monotonicity and the pointwise limit of $R_tu+Kt^2$.
\end{proof}

The classical theorem supplies a modulus $\epsilon$, not a rate claimed here to be linear in $\delta$. The cancellation and kernel argument of \cref{lem:chern-smoothing-integral} still give
\[
 \|R_tu-u\|_{L^1(\mathrm{d}V)}\leq Ct^2,
 \qquad \|\nabla R_tu\|_\infty\leq C/t.
\]
Indeed the averaging map is unchanged, $u$ is bounded, and the monotone lower bound just proved is all that argument uses.


\subsection{A linear error for a smooth background}

Let $\theta$ be a smooth big real $(1,1)$-form. Choose a strict potential $\rho$ with analytic poles and smooth resolved remainder as in \cref{thm:resolvable-regularization-input}, and put $\Omega=X\setminus\{\rho=-\infty\}$. The next estimate refines the curvature error for normalized potentials locally bounded on $\Omega$.

\Needspace{24\baselineskip}
\begin{lemma}\label{lem:smooth-background-regularization-rate}
Let $a$ be quasi-plurisubharmonic on $X$, with $\sup_Xa=0$ and
$\theta+\ddc a\geq0$, and suppose that $a$ is locally bounded on $\Omega$.
There are constants $A,K,C>0$ such that the averaging operators $R_t$
from \cref{lem:chern-kiselman-input} satisfy
\[
 R_ta+Kt^2\downarrow a,\qquad
 \|R_ta-a\|_{L^1(\mathrm{d}V)}\leq Ct^2.
\]
For small $\delta>0$ and $c>0$, the transform on $\Omega$
\[
 \Phi_{c,\delta}
 =\inf_{0<t\leq\delta}
 \{R_ta+K(t^2-\delta^2)+K(t-\delta)-c\log(t/\delta)\}
\]
extends canonically to $X$ and satisfies
\begin{equation}\label{eq:smooth-background-transform-rate}
 a-K(\delta^2+\delta)\leq\Phi_{c,\delta}\leq R_\delta a\leq0,
 \qquad \theta+\ddc\Phi_{c,\delta}\geq-(Ac+C\delta)\omega.
\end{equation}
At each point of $\Omega$ the infimum is attained at a positive radius.
On every compact $K_1\Subset\Omega$, for sufficiently small $t$,
$\|\nabla R_ta\|_{L^\infty(K_1)}\leq C_{K_1}/t$.
\end{lemma}

\begin{proof}
Apply the classical Chern-connection regularization expansion to the closed current $\ddc a$ and the smooth lower bound $-\theta$. It gives monotonicity and logarithmic
convexity: for a constant $K_0>0$, set
$\lambda(x,t)=t\partial_t(R_ta(x)+K_0t^2)$. This function is nonnegative and increasing in $t$.
We first justify a linear error for this smooth lower bound, rather than
infer a rate from the unspecified modulus in that theorem.

Choose the auxiliary curvature form to be $A\omega$, with $A$ large
enough that the Chern curvature plus $A\omega\otimes\operatorname{Id}$
is nonnegative on every decomposable tensor. Equivalently, after
absorbing the normalization constants, writing $R_{j\bar k\ell\bar m}$ for the Chern curvature coefficients and $\xi,\zeta$ for tangent vectors in a unitary frame, the coefficients satisfy
\[
 \sum R_{j\bar k\ell\bar m}\xi_j\bar\xi_k\zeta_\ell\bar\zeta_m
 \geq-A|\xi|^2|\zeta|^2.
\]
This choice is available by compactness. Pairing this inequality with the nonnegative eigenvalues of a positive complex Hessian bounds its curvature contraction by $-A$ times its trace, with no further error term.

Fix $x\in X$ and a complex radius variable $w\ne0$, and put $t=|w|$. Write $\xi\in T_x^{1,0}X$ and $\tau\in\mathbb C$ for a base--radius tangent vector. In a uniformly sized coordinate chart about $x$,
write $\beta=\ddc|z|^2$ and
choose a quadratic $q_{x,t}$ with
\[
 \ddc q_{x,t}=\theta(x)+Bt\beta,
\]
where $B$ is fixed and large enough that $a+q_{x,t}$ is plurisubharmonic
on a ball containing the averaging support and its nearby base points.
This follows from the smooth variation bound
$\|\theta(y)-\theta(x)\|\leq C|y-x|$ on that ball of radius $O(t)$.
For the Hessian calculation at $(x,w)$ the quadratic is held fixed.
Its coefficients and all needed derivatives are uniformly bounded.

We estimate the errors in the classical regularization expansion.
Normalized quasi-plurisubharmonic compactness bounds $\|a\|_1$.
On each fixed larger chart, add one fixed strictly plurisubharmonic
quadratic $Q$ to $a$. If
$\nu=\ddc(a+Q)\wedge\beta^{n-1}$, the local Lelong mass inequality and
a fixed larger-ball mass bound give
$\nu(B_t(x))\leq Ct^{2n-2}$. The coefficients of
$\partial\bar\partial(a+q_{x,t})$ are controlled in total variation
by those of $\ddc(a+Q)$ plus a uniformly bounded smooth form. Write $|\partial\bar\partial(a+q_{x,t})|$ for the sum of the total variations of its coefficient measures. The rescaled averaging kernels are supported in $B(x,Ct)$ and bounded by a constant times $t^{-2n}$. Their integrals against this measure are therefore bounded by
\[
 C t^{-2n}|\partial\bar\partial(a+q_{x,t})|(B(x,Ct))
 \leq C't^{-2}.
\]
Thus the terms with a coefficient $O(t^3)$ contribute
$O(t)$; the order-two expansion has an $O(t)$ remainder as well.
The mixed and vertical contributions have bounds
$C(|\xi||\tau|+|\tau|^2)$. The curvature term is bounded below by
$-A\lambda(x,t)|\xi|^2$, using the preceding curvature choice.
The radial term of the same expansion identifies this mass with $\lambda$ up to $O(t^2)$.
Subtracting the regularization of $q_{x,t}$ restores $-\theta(x)$;
its base Hessian differs from $\ddc q_{x,t}$ by $O(t)$, its mixed and
vertical Hessians are bounded, and its radial logarithmic derivative
is $O(t^2)$. The extra $Bt\beta$ also costs only $O(t)$.
For a real $(1,1)$-form $\alpha$, write $\alpha[v]^2$ for its associated Hermitian quadratic form evaluated on $v$. Consequently, for a fixed constant $C_1>0$, the full base--fiber Hessian satisfies
\[
 \ddc_{x,w}R_{|w|}a[\xi,\tau]^2
 \geq-\theta_x[\xi]^2-(A\lambda(x,t)+Ct)|\xi|^2
       -C_1(|\xi||\tau|+|\tau|^2).
\]
All constants are independent of $x,t,c,\delta$ and uniform in the
normalized family. In particular, setting $\xi=0$ in the same local
expansion gives a uniform lower bound for the fiber Hessian. Adding a
fixed multiple of $|w|^2$ gives its subharmonicity, hence the
logarithmic convexity and monotonicity above with a uniform $K_0$.

Increase $K\geq K_0$ and add $K(t^2+t)$ to the regularization. The
positive vertical Hessian of $K|w|$ absorbs the mixed term by
\[
 C_1|\xi||\tau|\leq C't|\xi|^2+\frac{K}{8t}|\tau|^2,
\]
and absorbs the remaining bounded vertical error for small $t$.
The logarithmic derivative of the corrected function is at least
$\lambda$. The capped Kiselman minimum argument in
\cref{lem:general-kiselman-modulus} therefore gives the last inequality
in \eqref{eq:smooth-background-transform-rate}, now with error $C\delta$.
On compact subsets of $\Omega$, local boundedness of $a$ makes the
logarithmic penalty diverge uniformly at zero, so that the minimum
occurs away from zero. The same argument handles minimizers at the
upper cap without differentiating their radius.

The other bounds in \eqref{eq:smooth-background-transform-rate} follow
from monotonicity and evaluation at $t=\delta$. In particular the
transform is bounded above near the analytic complement. The removable
singularity theorem extends it, preserving the exact smooth lower
Hessian bound by local freezing and passage to the limit. Its lower
bound by $a-K(\delta^2+\delta)$ also extends by canonical representatives.

The adjoint computation in \cref{lem:chern-smoothing-integral} gives
$R_t^*1=1+O(t^2)$. It needs only $a\in L^1$, since
\[
 \left|\int_X(R_ta-a)\,\mathrm{d}V\right|
 \leq Ct^2\|a\|_1.
\]
Together with $R_ta-a\geq-Kt^2$, this proves the stated $L^1$ estimate.
Finally the smoothing kernel near $K_1$ samples a fixed larger compact
subset of $\Omega$, where $a$ is bounded. The same kernel derivative
estimate as in \cref{lem:chern-smoothing-integral} gives $C_{K_1}/t$.
\end{proof}

\section{Pullback and descent}
\label{sec:psh-descent}

Pullback preserves the local positivity of a potential. Compact fibers force descent when the background has no vertical component; under a modification, a Jacobian factor accounts for the change in multiplier-ideal integrability.

\Needspace{18\baselineskip}
\subsection{Holomorphic disks}

A smooth background need not have local potentials, but its pullback to a disk does. Testing there retains the exact lower curvature bound under a holomorphic map.

\begin{lemma}\label{lem:smooth-disk-criterion}
Let $\theta$ be a smooth real $(1,1)$-form on a complex manifold $M$, and let $u$ be quasi-plurisubharmonic on $M$. Then $\theta+\ddc u\geq0$ if and only if the following condition holds. For every holomorphic map $f:\Delta\to M$ from the unit disk whose pullback $u\circ f$ is not identically $-\infty$, the function $u\circ f+a$ is subharmonic wherever $a$ is a smooth local potential for $f^*\theta$.
\end{lemma}
\begin{proof}
Suppose first that $\theta+\ddc u\geq0$. On a coordinate ball convolve with a smooth nonnegative radial kernel. The smooth functions $u_\epsilon$ satisfy
\[
 \ddc u_\epsilon\geq-\theta_\epsilon,
 \qquad\theta_\epsilon\longrightarrow\theta\quad\text{smoothly}.
\]
Write $u=p-r$, with $p$ plurisubharmonic and $r$ smooth. The convolutions $p_\epsilon$ decrease to $p$. On a smaller disk where $u\circ f$ is genuine, $p_\epsilon\circ f$ therefore converges in $L^1_{\mathrm{loc}}$ to $p\circ f$, while $r_\epsilon\circ f$ converges smoothly. Pull back the smooth inequality and pass to distributions. This gives $\ddc(u\circ f)+f^*\theta\geq0$. Since $u\circ f=p\circ f-r\circ f$ is quasi-subharmonic, the distributional inequality makes $u\circ f+a$ subharmonic.

Conversely, test on affine complex lines in a coordinate ball. For any fixed nonzero constant direction $\xi$, local integrability of $u$ implies that almost every parallel line has a genuine restriction. Fubini's theorem applied to nonnegative smooth test functions gives positivity of the scalar distribution obtained by evaluating the Hermitian coefficient matrix of $\theta+\ddc u$ on $\xi$. Positivity for a countable dense set of constant directions proves positivity of the Hermitian matrix of coefficient measures, hence the current inequality.
\end{proof}

Consequently, holomorphic pullback preserves the background inequality wherever the pulled-back potential is genuine: apply the disk criterion to compositions of maps. We also use descent of subharmonicity under a finite disk map $\rho(t)=t^k$, $k\geq1$. If $h\circ\rho$ is subharmonic, local inverse branches make $h$ subharmonic off zero. It is bounded above near zero and has the canonical extension value there, since
\[
 h(0)=(h\circ\rho)(0)
      =\limsup_{t\to0,\ t\ne0}h(t^k)
      =\limsup_{z\to0,\ z\ne0}h(z).
\]
The removable-singularity theorem makes $h$ subharmonic at zero as well.

\Needspace{17\baselineskip}
\subsection{Compact fibers}

The connected components of a submersion fiber determine a finite covering of the base. Potentials first descend to that covering.

\begin{lemma}\label{lem:submersion-component-cover}\index{Descent!cover of fiber components}\index{Submersion!cover of fiber components}
Let $p:M\to Y$ be a surjective holomorphic submersion between connected compact complex manifolds. There is a factorization
\[
 M\xrightarrow{\pi}Y'\xrightarrow{r}Y,\qquad p=r\circ\pi,
\]
where $Y'$ is a connected compact complex manifold, $r$ is a finite unramified holomorphic covering and $\pi$ is a surjective holomorphic submersion with connected fibers. The points of $r^{-1}(y)$ correspond to the connected components of $p^{-1}(y)$, which are the fibers of $\pi$ over those points.
\end{lemma}
\begin{proof}
Fix a smooth Hermitian metric on $M$. The orthogonal complement of $\ker\mathrm dp$ gives a smooth horizontal lift of tangent vectors on $Y$. Around $y_0\in Y$ choose a coordinate ball $B$ centered at $y_0$, with compact closure inside a larger coordinate ball. For $y\in B$, lift the straight segment from $y_0$ to $y$. The ordinary differential equation for horizontal transport has a solution along the whole segment: its base remains in the compact closure of $B$, whose inverse image is compact. Smooth dependence on the starting point and on $y$, together with transport along the reversed segment, gives a smooth trivialization
\[
 p^{-1}(B)\simeq B\times p^{-1}(y_0)
\]
over $B$.

The compact smooth manifold $p^{-1}(y_0)$ has finitely many connected components. The trivialization partitions $p^{-1}(B)$ into finitely many open and closed pieces, each with one connected component in every fiber. On overlaps, two such local partitions identify the same fiber components; their labels differ by locally constant permutations. Glue one copy of each base ball for every component using these identifications. Equivalently, take the quotient of $M$ that identifies points in the same connected component of a fiber. The local partitions give this quotient a Hausdorff covering topology: distinct points over one base point are separated by different sheets, and points over different base points are separated downstairs.

Call the quotient $Y'$ and its projection $r$. Base coordinates on each sheet give holomorphic charts, with holomorphic transition maps, so $r$ is a local biholomorphism. It is a finite covering and $Y'$ is compact, since a finite cover of a compact base is compact. The quotient map $\pi$ is locally $p$ restricted to one of the open pieces. It is therefore a holomorphic submersion, with the asserted connected fibers. Surjectivity of $\pi$ and connectedness of $M$ make $Y'$ connected.
\end{proof}

\Needspace{17\baselineskip}
\begin{proposition}\label{prop:psh-descent-submersion}\index{Descent!submersion}
Let $\pi:Y\to X$ be a surjective holomorphic submersion between connected compact complex manifolds, with connected fibers. Let $\theta$ be a smooth real $(1,1)$-form on $X$. Pullback gives a bijection
\[
 \PSH(X,\theta)\longrightarrow\PSH(Y,\pi^*\theta).
\]
Every member on the right is constant on each fiber, with the value $-\infty$ allowed on an entire fiber.
\end{proposition}
\begin{proof}
In a coordinate ball $V\subset X$, choose a smooth function $\rho$ with $\ddc\rho\geq\theta$. If $u\in\PSH(Y,\pi^*\theta)$, then $u+\rho\circ\pi$ is plurisubharmonic on $\pi^{-1}(V)$. Its restriction to a fiber is plurisubharmonic or identically $-\infty$. Since a fiber is compact and connected, the maximum principle makes it constant in either case. Thus $u=v\circ\pi$ for a uniquely defined function $v$ on $X$.

A holomorphic submersion admits local holomorphic sections. Composing $u+\rho\circ\pi$ with such a section shows that $v+\rho$ is plurisubharmonic or identically $-\infty$ on a smaller ball. The latter alternative on a nonempty open ball would make $u=-\infty$ on its open inverse image, contradicting genuineness. Hence $v$ is quasi-plurisubharmonic. Distributional differentiation in local product coordinates gives
\[
 \pi^*(\theta+\ddc v)=\pi^*\theta+\ddc u\geq0.
\]
Positivity of a current of this pulled-back form implies positivity on the base: test in product coordinates and integrate a nonnegative compactly supported volume form of integral one in the fiber variables. It follows that $v\in\PSH(X,\theta)$.

Conversely, \cref{lem:smooth-disk-criterion} shows that pullback of a $\theta$-plurisubharmonic function by a submersion is $\pi^*\theta$-plurisubharmonic. It is genuine by surjectivity and local product coordinates. Uniqueness follows from surjectivity.
\end{proof}

\Needspace{11\baselineskip}
\subsection{Modifications}

\begin{proposition}\label{prop:psh-descent-modification}\index{Descent!modification}
Let $\pi:Y\to X$ be a modification of connected compact complex manifolds, and let $\theta$ be a smooth real $(1,1)$-form on $X$. Pullback gives a bijection
\[
 \PSH(X,\theta)\longrightarrow\PSH(Y,\pi^*\theta).
\]
\end{proposition}
\begin{proof}
Let $S\subset X$ be a proper analytic set over whose complement $\pi$ is an isomorphism, and put $D=\pi^{-1}(S)$. Pullback preserves the curvature inequality by \cref{lem:smooth-disk-criterion} and is genuine on the isomorphism locus. It is injective because $\pi$ is surjective.

To prove surjectivity, let $b\in\PSH(Y,\pi^*\theta)$ and define $u$ off $S$ by the inverse map. On a coordinate neighborhood choose a smooth $\rho$ with $\ddc\rho\geq\theta$. Then $b+\rho\circ\pi$ is plurisubharmonic. Properness and its local upper bounds give a local upper bound for $u+\rho$ near $S$, so the removable-singularity theorem extends $u+\rho$ plurisubharmonically across $S$. The resulting quasi-plurisubharmonic extensions of $u$ agree almost everywhere on overlaps and hence glue as canonical representatives.

The current $\theta+\ddc u$ is positive off $S$. Locally it is the sum of the positive current $\ddc(u+\rho)$ and the smooth form $\theta-\ddc\rho$. The latter has zero coefficient-measure mass on $S$, so the restriction of $\theta+\ddc u$ to $S$ is positive as well. Thus $u\in\PSH(X,\theta)$. The plurisubharmonic functions $b+\rho\circ\pi$ and $(u+\rho)\circ\pi$ agree off $D$, hence almost everywhere and everywhere. Therefore $b=u\circ\pi$.
\end{proof}

A modification preserves $L^1$ convergence within a family with one lower curvature bound. Compactness controls the pullbacks near the exceptional locus, where a direct change-of-variables estimate would involve a vanishing Jacobian.

\begin{lemma}\label{lem:modification-lone-pullback}
Let $\pi:Y\to X$ be a modification between connected complex manifolds, and let $\theta$ be a smooth real $(1,1)$-form on $X$. If $v_j\to v$ in $L^1_{\mathrm{loc}}(X)$ and all $v_j,v$ belong to $\PSH(X,\theta)$, then
\[
 v_j\circ\pi\longrightarrow v\circ\pi
       \quad\hbox{in }L^1_{\mathrm{loc}}(Y).
\]
In particular, on compact manifolds this is convergence in $L^1$.
\end{lemma}
\begin{proof}
\Cref{lem:qpsh-hartogs} bounds $v_j$ uniformly above on relatively compact subsets of $X$. The pullbacks are therefore locally uniformly bounded above and have the common smooth lower curvature bound $-\pi^*\theta$. After a smooth local correction, plurisubharmonic compactness gives a convergent subsequence on each connected coordinate region in $Y$, or convergence to $-\infty$ locally uniformly.

Every such region meets the complement of the exceptional locus. There $\pi$ is a biholomorphism, and change of variables gives $L^1_{\mathrm{loc}}$ convergence to $v\circ\pi$. Thus collapse is excluded. Each subsequential limit agrees with $v\circ\pi$ off the exceptional locus, hence almost everywhere, since that analytic locus has smooth volume zero. The limit is unique. Applying this argument to every subsequence and a countable coordinate cover proves convergence of the whole sequence in $L^1_{\mathrm{loc}}(Y)$.
\end{proof}

Capacity convergence is also preserved by a modification. The comparison takes place away from small neighborhoods of the exceptional sets; a uniform comparison of capacities on the whole manifolds is unnecessary.

\Needspace{10\baselineskip}
\begin{lemma}\label{lem:modification-capacity-pullback}
Let $\pi:Y\to X$ be a modification between connected compact complex manifolds, equipped with Hermitian forms $\omega_Y,\omega_X$. Write $\operatorname{Cap}_Y=\operatorname{Cap}_{\omega_Y}$ and $\operatorname{Cap}_X=\operatorname{Cap}_{\omega_X}$. For quasi-plurisubharmonic functions $v_j,v$ on $X$,
\[
 v_j\longrightarrow v\text{ in capacity on }X
 \quad\Longleftrightarrow\quad
 v_j\circ\pi\longrightarrow v\circ\pi\text{ in capacity on }Y.
\]
\end{lemma}
\begin{proof}
Choose a proper analytic set $S\subset X$ such that $\pi$ is an isomorphism outside $D=\pi^{-1}(S)$. Both sets have capacity zero. For $\epsilon>0$, choose an open neighborhood $V$ of $D$ of capacity less than $\epsilon$. Properness gives a neighborhood $U_0$ of $S$ with $\pi^{-1}(U_0)\subset V$. Shrink it to an open $U$ with $\operatorname{Cap}_X(U)<\epsilon$. The inverse map is biholomorphic near the compact set $X\setminus U$. A finite cover by nested coordinate charts and the local-global capacity comparison of \cref{sec:hermitian-capacity} give a constant $C_\epsilon$ such that, for every Borel set $E\subset X$,
\[
 \begin{split}
 \operatorname{Cap}_Y(\pi^{-1}E)&\leq\epsilon+C_\epsilon\operatorname{Cap}_X(E),\\
 \operatorname{Cap}_X(E)&\leq\epsilon+C_\epsilon\operatorname{Cap}_Y(\pi^{-1}E).
 \end{split}
\]
Apply these inequalities to $E=\{|v_j-v|>\delta\}$ for each fixed $\delta>0$, omitting the pluripolar pole sets. First let $j\to\infty$, then $\epsilon\downarrow0$.
\end{proof}

\Needspace{13\baselineskip}
\subsection{Multiplier ideals under modifications}
\label{sec:psh-birational-ideals}

Let $\pi:Y\to X$ be a modification of connected compact complex manifolds of dimension $n$. There is an analytic subset $S\subset X$ of codimension at least two such that $\pi$ is an isomorphism over $X\setminus S$. Let $K_X$ and $K_Y$ be the canonical holomorphic line bundles. Set $D=\pi^{-1}(S)$ and
\[
 R_\pi=K_Y\otimes\pi^*K_X^{-1}.
\]
The determinant of the differential is a canonical holomorphic section $J_\pi$ of $R_\pi$. In local coordinates its coefficient is the complex Jacobian determinant. Its zeros are contained in $D$. These standard facts about modifications of smooth manifolds include the case $n=1$, when the modification is an isomorphism.

\Needspace{20\baselineskip}
\begin{lemma}\label{lem:adjoint-ideal-transform}\index{Multiplier ideal!birational transform}
Let $a$ be quasi-plurisubharmonic on $X$. Under the natural identification away from $S$,
\[
 \pi_*\bigl(R_\pi\otimes\mathcal I(a\circ\pi)\bigr)
       =\mathcal I(a).
\]
More precisely, for any holomorphic line bundle $M$ on $X$, the map
\[
 s\longmapsto J_\pi\,\pi^*s
\]
identifies
\[
 H^0\bigl(X,M\otimes\mathcal I(a)\bigr)
 \simeq H^0\bigl(Y,\pi^*M\otimes R_\pi
                         \otimes\mathcal I(a\circ\pi)\bigr).
\]
\end{lemma}
\begin{proof}
The assertion is local on $X$. In coordinates and local frames the proposed map sends a holomorphic coefficient $f$ to $(f\circ\pi)\operatorname{Jac}\pi$. Change of variables gives the local integrability equivalence between
\[
 |f|^2\mathrm{e}^{-a}
 \quad\hbox{and}\quad
 |f\circ\pi|^2|\operatorname{Jac}\pi|^2\mathrm{e}^{-a\circ\pi}.
\]
For the implication from integrability upstairs to integrability near a downstairs point, take a relatively compact smaller neighborhood. Its inverse image has compact closure by properness, so finitely many integrability charts suffice. Smooth positive coordinate densities are comparable on these smaller charts. Analytic exceptional sets have smooth volume zero and do not change the integrals.

It remains to justify that every section upstairs has the required form. Divide it by $J_\pi$ off $D$ and use the inverse of $\pi$ there. This gives a holomorphic section of $M$ off $S$. Hartogs extension across the codimension-at-least-two analytic set $S$ supplies a holomorphic section $s$ on $X$. Its image $J_\pi\pi^*s$ agrees with the given section off $D$, hence everywhere by the identity theorem. The integrability equivalence then places $s$ in the required multiplier ideal. The same argument on each open set proves the sheaf identity and the global identification.
\end{proof}

In particular, the formula does not identify $\pi^*\mathcal I(a)$ with $\mathcal I(a\circ\pi)$. The Jacobian is part of the holomorphic coefficient upstairs, rather than a factor that may be deleted from the integrability test.

We also need invariance of the ideal \emph{order}, which asks for all positive weights. We give the local argument to avoid assuming that a modification of a small open subset of $Y$ extends over $X$.

\Needspace{10\baselineskip}
\begin{lemma}\label{lem:birational-ideal-order}
For quasi-plurisubharmonic $a,b$ on $X$,
\[
 a\preceq_{\mathcal I}b
 \quad\Longleftrightarrow\quad
 a\circ\pi\preceq_{\mathcal I}b\circ\pi.
\]
The assertion is local on $X$ and also holds for a proper bimeromorphic map between noncompact smooth complex manifolds.
\end{lemma}
\begin{proof}
The reverse implication follows at every positive weight from the sheaf identity in \cref{lem:adjoint-ideal-transform}: tensor the upstairs ideal inclusion with $R_\pi$ and apply the left-exact direct image functor. The proof of that sheaf identity is local on the target; properness makes each fiber compact and supplies the required finite integrability cover. Compactness of the entire manifolds is unnecessary.

For the forward implication, fix a coordinate neighborhood $W\subset Y$, a smooth local modification $\sigma:Z\to W$, and a prime divisor $E$ on $Z$. Shrinking $W$ as necessary near the center under consideration, its image is contained in a coordinate neighborhood $U\subset X$. Put $F=\pi\circ\sigma$ and
\[
 A_F(E)=1+\operatorname{ord}_E\operatorname{Jac}F.
\]
The map $F$ is holomorphic and is one-to-one and biholomorphic off a proper analytic set. It need not be a proper modification of $U$.

For a positive integer $m$, let $\nu_E$ denote the generic Lelong number along $E$, and let $\operatorname{ord}_E$ denote the generic order of an ideal there. The proof of \cref{lem:bergman-divisorial-error} applies to this map on relatively compact coordinate patches and gives
\[
 \max\left\{\nu_E(a\circ F)-\frac{A_F(E)}m,0\right\}
 \leq \frac1m\operatorname{ord}_E\bigl(\mathcal I(ma)\mathcal O_Z\bigr)
 \leq\nu_E(a\circ F).
\]
For clarity, the upper inequality pulls back the local Bergman lower bound on $U$. For the lower one, a germ $f\in\mathcal I(ma)$ pulls back to an integrable coefficient $(f\circ F)\operatorname{Jac}F$: on a relatively compact upstairs patch the change-of-variables integral is bounded by the integral over its downstairs image. Generic divisorial removal and the one-variable integral give
\[
 \operatorname{ord}_E(f\circ F)+A_F(E)>m\nu_E(a\circ F).
\]
Only local injectivity off the exceptional set and the compact patch are needed for this estimate, not properness of $F$ over $U$. The same estimates hold for $b$.

The hypothesis implies $\mathcal I(ma)\subseteq\mathcal I(mb)$ for every $m>0$. Orders of the pulled-back ideals therefore satisfy the reverse inequality. Letting $m\to\infty$ in the two estimates proves
\[
 \nu_E(a\circ F)\geq\nu_E(b\circ F).
\]
This holds for every local divisor over $Y$. Blowing up a point on any local modification recovers its pointwise Lelong number; in dimension one the point already is a divisor. The exact local criterion in \cref{thm:local-i-inputs} now proves the upstairs ideal order. No extension of $\sigma$ outside its local domain was used.
\end{proof}

\Needspace{13\baselineskip}
\begin{lemma}\label{lem:ideal-equivalence-modification}
For a proper bimeromorphic map $\rho:Z\to Y$ between smooth complex manifolds and genuine quasi-plurisubharmonic functions $a,b$ on $Y$,
\[
 a\sim_{\mathcal I}b
       \quad\Longleftrightarrow\quad
 a\circ\rho\sim_{\mathcal I}b\circ\rho.
\]
The statement is local on $Y$.
\end{lemma}
\begin{proof}
Apply the local order equivalence of \cref{lem:birational-ideal-order} to both $(a,b)$ and $(b,a)$. The two order inequalities are exactly equality of the multiplier ideals at every positive weight.
\end{proof}

\Needspace{20\baselineskip}
\subsection{Restriction of analytic approximants}

A quasi-plurisubharmonic function with zero generic Lelong number along a submanifold may still restrict to the function identically equal to $-\infty$ (\cref{sec:analytic-types}). Analytic regularization removes this obstruction: the restrictions of the analytic approximants are quasi-plurisubharmonic, and a fixed lower curvature bound controls their generic Lelong numbers along divisors on every fixed modification.

\begin{corollary}\label{cor:restricted-analytic-divisorial-bound}
Let $\theta$ be a smooth real $(1,1)$-form and $u\in\PSH(X,\theta)$, and let $u_m$, with curvature errors $\epsilon_m$, be the approximants of \cref{cor:nonclosed-analytic-approximation}. Let $Y\subset X$ be a connected compact complex submanifold along which the generic Lelong number of $u$ vanishes:
\[
 \nu(u,Y)=\inf_{y\in Y}\nu(u,y)=0.
\]
Then each restriction $u_m|_Y$ is quasi-plurisubharmonic on $Y$; in particular, it is not identically $-\infty$. If $\pi\colon Z\to Y$ is a modification with smooth source and $E\subset Z$ is a prime divisor, then
\[
 \sup_m\nu\bigl((u_m|_Y)\circ\pi,E\bigr)<\infty.
\]
\end{corollary}
\begin{proof}
For each $m$, choose $y\in Y$ with $\nu(u,y)<1/m$. Skoda's local integrability criterion gives $\mathcal I(mu)_y=\mathcal O_{X,y}$, so the image ideal $\mathcal I(mu)\mathcal O_Y$ does not vanish identically. Since $u_m$ has analytic type $(1/m,\mathcal I(mu))$, \cref{prop:analytic-type-calculus} shows that $u_m|_Y$ is quasi-plurisubharmonic on $Y$. Its pullback to $Z$ is quasi-plurisubharmonic as well, since $\pi$ is biholomorphic off a proper analytic subset.

\Needspace{6\baselineskip}
Restricting and pulling back the curvature inequality of \cref{cor:nonclosed-analytic-approximation}, which is allowed by \cref{lem:smooth-disk-criterion}, gives
\[
 \pi^*(\theta|_Y)+\ddc\bigl((u_m|_Y)\circ\pi\bigr)
 \geq-\epsilon_m\pi^*(\omega|_Y).
\]
Since the errors $\epsilon_m$ are bounded, a sufficiently large multiple of a single Hermitian form on $Z$ dominates $\pi^*(\theta|_Y)+\epsilon_m\pi^*(\omega|_Y)$ for all $m$. \Cref{lem:gauduchon-uniform-lelong}, applied on $Z$ in this fixed background, therefore bounds the generic Lelong numbers along $E$ uniformly in $m$.
\end{proof}


\section{Convergence against a non-pluripolar measure}
\label{sec:psh-fixed-measure}

A common bound on the potentials allows $L^1$ convergence for smooth volume to be transferred to a fixed finite measure that does not charge pluripolar sets. Bounded exponentials extend this statement to upper-bounded families.

Cegrell's fixed-measure convergence lemma says that a uniformly bounded sequence of plurisubharmonic functions converging in $L^1_{\mathrm{loc}}$ converges in $L^1(\nu)$ on compact subsets for each finite non-pluripolar measure $\nu$ there. The usual formulation allowing extraction gives the same conclusion: apply it to every subsequence, extract almost-everywhere convergence for smooth volume, and use the subsequence criterion for convergence. This is the classical fixed-measure argument used in \cite[Section 11.1.2]{GZ17}. We retain this local lemma as a standard input.

\Needspace{15\baselineskip}
\begin{lemma}\label{lem:fixed-npp-integrals}
Let $\nu$ be a finite non-pluripolar measure on $X$. If uniformly bounded quasi-plurisubharmonic functions $a_j,a$ have a common smooth lower curvature bound and $a_j\to a$ in $L^1(X)$, then
\[
 \int_X|a_j-a|\,\mathrm{d}\nu\longrightarrow0.
\]
If instead $u_j,u\in\PSH(X,\theta)$ have a common upper bound and $u_j\to u$ in $L^1(X)$, then, for every $b>0$,
\[
 \int_X|\mathrm{e}^{bu_j}-\mathrm{e}^{bu}|\,\mathrm{d}\nu\longrightarrow0.
\]
\end{lemma}
\begin{proof}
For the first statement add a fixed smooth strictly plurisubharmonic correction on each of finitely many coordinate balls. The resulting functions are uniformly bounded and plurisubharmonic. Apply the local input on smaller compact balls covering $X$ and sum the integrals; the corrections cancel in the difference.

For the second statement choose $A\omega\geq\theta$. If $u_j\leq C$, the convex chain rule gives
$\ddc \mathrm{e}^{bu_j}\geq-bA\mathrm{e}^{bC}\omega$, also for unbounded $u_j$ by bounded truncation. Thus $\mathrm{e}^{bu_j}$ are uniformly bounded quasi-plurisubharmonic functions in one lower curvature class. Their $L^1(\mathrm{d}V)$ convergence follows from the Lipschitz bound $b\mathrm{e}^{bC}$ for the exponential on $(-\infty,C]$. Apply the first statement. We assign the value zero to the exponential at a pole.
\end{proof}

\begin{proposition}\label{prop:singular-no-uniform-limit}
Let $\theta$ be smooth and let $\varphi\in\PSH(X,\theta)$. Suppose $\varphi$ is unbounded below and $u_\lambda\sim\varphi$ for every $\lambda>0$. Then $\lambda u_\lambda$ cannot converge uniformly on $X$ to a finite constant as $\lambda\downarrow0$. Its distance to any finite constant is also infinite in $L^\infty(\mathrm{d}V)$ for each fixed $\lambda>0$.
\end{proposition}
\begin{proof}
Bounded equivalence preserves unboundedness below. Every nonempty strict sublevel set of the upper semicontinuous function $u_\lambda$ is open and has positive smooth volume. Thus $u_\lambda$ is essentially unbounded below, as is $\lambda u_\lambda$. The stated conclusions follow.
\end{proof}

A vanishing exponent gives convergence to one even when the normalized potentials vary.

\Needspace{17\baselineskip}
\begin{lemma}\label{lem:normalized-exponential-lp}\label{lem:normalized-exponential-npp}
Let $v_j\in\PSH(X,\theta)$ satisfy $\sup_Xv_j=0$, and let $\lambda_j>0$ decrease to zero. For every finite non-pluripolar measure $\mu$,
\[
 \int_X\left|\mathrm{e}^{\lambda_jv_j}-1\right|\,\mathrm{d}\mu\longrightarrow0.
\]
If $f\in L^p(\mathrm{d}V)$ for some $p>1$, then more precisely
\[
 \int_X\left|\mathrm{e}^{\lambda_jv_j}-1\right|\,|f|\,\mathrm{d}V
 \leq C_p\lambda_j\|f\|_{L^p(\mathrm{d}V)}.
\]
The constant depends on the fixed background, $p$ and $\mathrm{d}V$.
\end{lemma}
\begin{proof}
Choose a Hermitian form $\Omega\geq\theta$. By \cref{lem:global-capacity-estimates}, for every $T>1$,
\[
 \operatorname{Cap}_\Omega(\{v_j<-T\})\leq C/T
\]
uniformly in $j$. Since $v_j\leq0$, splitting at this sublevel gives
\[
 \int_X(1-\mathrm{e}^{\lambda_jv_j})\,\mathrm{d}\mu
 \leq\lambda_jT\mu(X)+\mu(\{v_j<-T\}).
\]
\Cref{lem:npp-uniform-capacity} makes the second term uniformly small when $T$ is large. For this fixed $T$, the first term tends to zero as $j\to\infty$. This proves the general assertion. The value of the exponential is zero at a pole; pole sets have zero $\mu$-measure.

For the quantitative assertion, \cref{lem:uniform-exponential}, applied in background $\Omega$, gives a uniform $L^q$ bound for $-v_j$, where $q=p/(p-1)$. H\"older's inequality and
\[
 0\leq1-\mathrm{e}^{\lambda_jv_j}\leq\lambda_j|v_j|
\]
give the claimed estimate for every $j$.
\end{proof}

A moving bounded-product measure can also be used if its potentials converge in capacity. The proof estimates the torsion terms in integration by parts.

\begin{lemma}\label{lem:moving-bounded-integrals}
Let $a_j,a,w_j,w\in\PSH(X,\omega)$ be uniformly bounded. If $a_j\to a$ in $L^1(X)$ and $w_j\to w$ in capacity, then, for $0\leq k\leq n$,
\[
 \int_X|a_j-a|\,\omega_{w_j}^{k}\wedge\omega^{n-k}
 \longrightarrow0.
\]
The same conclusion holds when the $a_j$ have any fixed smooth lower curvature bound.
\end{lemma}
\begin{proof}
By the subsequence criterion, we may first extract $a_j\to a$ almost everywhere for smooth volume. The upper tail envelopes
$A_j=(\sup_{\ell\geq j}a_\ell)^*$ decrease to $a$: the local Hartogs lemma gives the upper inequality for the limit, and almost-everywhere convergence gives the reverse one. Bounded monotone capacity convergence therefore implies
\[
 \alpha_j:=\max(a_j,a)\longrightarrow a\quad\text{in capacity}.
\]
The part $\alpha_j-a$ of $|a_j-a|=(\alpha_j-a)+(\alpha_j-a_j)$ tends to zero against the moving products by \cref{lem:bounded-mixed-capacity}.

Put $h_j=\alpha_j-a_j\geq0$. We compare its integrals for two indices $j,\ell$. The case $k=0$ is the assumed $L^1$ convergence. For $k\geq1$, put $z=w_j-w_\ell$ and
\[
 T=\sum_{s=0}^{k-1}\omega_{w_j}^{s}\wedge
       \omega_{w_\ell}^{k-1-s}\wedge\omega^{n-k}.
\]
The sum includes $s=0$. Factorization and compact Stokes give
\begin{align*}
 &\int_Xh_j\bigl(\omega_{w_j}^{k}-\omega_{w_\ell}^{k}\bigr)
                      \wedge\omega^{n-k}\\
 &\hspace{20mm}=\int_Xh_j\ddc z\wedge T
 =\int_Xz\,\ddc(h_jT).
\end{align*}
Here the full product rule is
\begin{equation}\label{eq:moving-measure-torsion}
 \ddc(h_jT)=\ddc h_j\wedge T+\mathrm{d}h_j\wedge\dc T
             -\dc h_j\wedge \mathrm{d}T+h_j\ddc T.
\end{equation}
The last term is retained. It need not vanish on a Hermitian manifold.

We estimate the four terms. Since $\ddc h_j=\omega_{\alpha_j}-\omega_{a_j}$, the absolute contribution of the first is at most
\[
 \int_X|z|(\omega_{\alpha_j}+\omega_{a_j})\wedge T.
\]
It tends to zero uniformly as $j,\ell\to\infty$ by \cref{lem:bounded-mixed-capacity}, since $w_j-w_\ell\to0$ in capacity. Differentiating each factor shows
$\mathrm{d}T=\mathrm{d}\omega\wedge R$ and $\dc T=\dc\omega\wedge R$, where $R$ is a positive sum of mixed products of degree $(n-2,n-2)$. Moreover $\ddc T$ is a sum of terms containing $\ddc\omega$ and mixed products of degree $n-2$, or $\mathrm{d}\omega\wedge\dc\omega$ and mixed products of degree $n-3$. Terms with negative degree are absent. These smooth torsion factors are bounded in absolute value by fixed multiples of $\omega^2$ and $\omega^3$. As $h_j$ is uniformly bounded, the contribution of $h_j\ddc T$ also tends uniformly to zero by the same capacity estimate.

For the two gradient terms, the pointwise Cauchy--Schwarz inequality for a positive mixed product gives, for example,
\[
 \left|\int_Xz\,\mathrm{d}h_j\wedge\dc\omega\wedge R\right|^2
 \leq C\left(\int_X|z|\,\mathrm{d}h_j\wedge\dc h_j\wedge\omega\wedge R\right)
          \left(\int_X|z|\,\omega^2\wedge R\right).
\]
The second factor tends uniformly to zero. The first is uniformly bounded. For completeness, if $p$ is any of the uniformly bounded potentials and $S$ any positive mixed product of degree $n-1$ occurring here, then
\[
 \int_X\mathrm{d}p\wedge\dc p\wedge S
 =\frac12\int_Xp^2\ddc S-\int_Xp\,\ddc p\wedge S.
\]
Expand $\ddc S$ with all its torsion terms as above, and write $\ddc p=\omega_p-\omega$. Boundedness of $p$ and \cref{lem:bounded-mixed-capacity} bound both integrals uniformly. Finally,
$\mathrm{d}h_j\wedge\dc h_j\leq2\mathrm{d}\alpha_j\wedge\dc\alpha_j+2\mathrm{d}a_j\wedge\dc a_j$
as positive quadratic forms. This proves the asserted bound. The other gradient term has the same estimate. In dimension one $T=1$ and all torsion terms in \eqref{eq:moving-measure-torsion} vanish.

These bounded integration-by-parts and energy identities may first be written for smooth decreasing regularizations in the exact Hermitian background. The local Bedford--Taylor chain rule, its gradient products and quasi-continuity give the displayed identities for bounded potentials. Equivalently they follow from the bounded current calculus of \cref{sec:psh-bounded-products}. Only the common bound on the potentials enters the estimates.

We have shown that the difference of the two $h_j$ integrals tends uniformly to zero when $j,\ell$ are large. For a fixed sufficiently large $\ell$, \cref{lem:fixed-npp-integrals} gives
\[
 \int_Xh_j\,\omega_{w_\ell}^{k}\wedge\omega^{n-k}\longrightarrow0.
\]
Then let $\ell\to\infty$. Together with the estimate for $\alpha_j-a$, this proves the result. If the $a_j$ have a larger fixed lower curvature bound, enlarge $\omega$ to $\Omega$ so that they are $\Omega$-plurisubharmonic. Apply the proved statement with $\Omega$; positivity bounds the original mixed measures by their $\Omega$ counterparts.
\end{proof}

\subsection{Smoothing the measure}
\label{subsec:psh-measure-smoothing}

Coordinate convolution gives smooth approximations of a finite measure. For a varying quasi-plurisubharmonic integrand, the submean inequality transfers the error to the original measure. This is the additional control needed beyond weak convergence of the measures themselves.

For a smooth background $\theta$ with $\PSH(X,\theta)\ne\varnothing$ and a non-pluripolar probability measure $\nu$, put
\[
 A_\theta(\nu)=\sup\left\{\int_X(-w)\,\mathrm{d}\nu:
 w\in\PSH(X,\theta),\ \sup_Xw=0\right\}.
\]
This number may be infinite. Bounded Bedford--Taylor measures satisfy the following estimate.

\begin{lemma}\label{lem:bounded-carrier-integrability}
Let $\omega$ be Hermitian, let $\psi\in\PSH(X,\omega)$ be bounded, and put $\nu=(\omega+\ddc\psi)^n$. There is a finite constant $C$, depending on the fixed backgrounds and $\|\psi\|_\infty$, such that
\[
 \int_X(-w)\,\mathrm{d}\nu\leq C
 \qquad(w\in\PSH(X,\theta),\ \sup_Xw=0).
\]
\end{lemma}
\begin{proof}
On finitely many nested coordinate balls $B'_a\Subset B_a$ covering $X$, choose fixed smooth corrections $q_a,r_a$ with $\ddc q_a\geq\omega$ and $\ddc r_a\geq\theta$, and shift $r_a$ to be nonpositive. Then $\psi+q_a$ is bounded plurisubharmonic, $w+r_a$ is negative plurisubharmonic, and positivity of bounded products gives
\[
 \int_{B'_a}(-w)\,\mathrm{d}\nu
 \leq\int_{B'_a}(-w-r_a)\bigl(\ddc(\psi+q_a)\bigr)^n
 \leq C_a\bigl(1+\|w\|_{L^1(B_a)}\bigr).
\]
The last step is the weighted Chern--Levine--Nirenberg estimate in \cref{sec:psh-local-capacity}. Sum over the cover and use the normalized $L^1$ bound of \cref{thm:normalized-compactness}.
\end{proof}

Choose finitely many coordinate balls $B_a$, smooth functions $\eta_a\geq0$ with compact support in these balls and $\sum_a\eta_a=1$, and inner neighborhoods of their supports whose closures lie in $B_a$. Let $\chi\geq0$ be a smooth radial convolution kernel on the unit ball of $\mathbb C^n$, of integral one. Fix $\epsilon_j\downarrow0$ small enough for every convolution below to stay in the corresponding coordinate ball. Denote convolution of a function in that chart by $S_{a,j}$. For a Borel test function $F$ and a probability measure $\nu$, define a smooth probability measure $\nu_j$ by convolving each coordinate measure $\eta_a\nu$ and adding its extensions by zero. Its defining dual identity is
\begin{equation}\label{eq:measure-smoothing-duality}
 \int_X F\,\mathrm{d}\nu_j
 =\sum_a\int_X\eta_a S_{a,j}F\,\mathrm{d}\nu.
\end{equation}
The kernel is even, so the same convolution occurs on the right. This identity first holds for nonnegative Borel functions, and then for integrable signed functions. In particular $\nu_j\rightharpoonup\nu$.

\Needspace{20\baselineskip}
\begin{lemma}\label{lem:measure-smoothing-integrals}
If $\nu$ is non-pluripolar and $A_\theta(\nu)<\infty$, then
\[
 \sup_j A_\theta(\nu_j)<\infty.
\]
More generally, let $z_j,z$ be quasi-plurisubharmonic functions with a common upper bound and a common smooth lower curvature bound. Suppose they converge both in $L^1(X)$ and in $L^1(\nu)$. Then
\begin{equation}\label{eq:convolution-moving-integrals}
 \int_X|z_j-z|\,\mathrm{d}\nu_j\longrightarrow0.
\end{equation}
Consequently, for every $\lambda>0$,
\[
 \mathrm{e}^{\lambda z_j}\nu_j\rightharpoonup \mathrm{e}^{\lambda z}\nu.
\]
\end{lemma}
\Needspace{9\baselineskip}
\begin{proof}
On each coordinate ball choose a fixed smooth function $q_a$ whose Levi form dominates $\theta$. If $w\in\PSH(X,\theta)$ and $\sup_Xw=0$, then $w+q_a$ is plurisubharmonic. The submean inequality gives
\[
 -S_{a,j}w\leq-w+S_{a,j}q_a-q_a\leq-w+C.
\]
The constant is independent of $w$ and $j$. Integrating \eqref{eq:measure-smoothing-duality} proves $A_\theta(\nu_j)\leq A_\theta(\nu)+C$.

For the second assertion start with any subsequence, and extract so that $z_j\to z$ almost everywhere for smooth volume. The upper envelopes
\[
 H_j=\left(\sup_{k\geq j}z_k\right)^*
\]
decrease to $z$ by \cref{thm:normalized-compactness} and this almost-everywhere convergence. Enlarge the fixed corrections $q_a$ so that $H_j+q_a,z_j+q_a,z+q_a$ are all plurisubharmonic. Since $H_j\geq z_j,z$, the dual identity and the submean inequality imply
\[
 \begin{split}
 \int_X|z_j-z|\,\mathrm{d}\nu_j
 &\leq\sum_a\int_X\eta_a S_{a,j}(2H_j-z_j-z)\,\mathrm{d}\nu\\
 &\leq2\sum_a\int_X\eta_a
       \bigl(S_{a,j}(H_j+q_a)-(z+q_a)\bigr)\,\mathrm{d}\nu
       +\int_X(z-z_j)\,\mathrm{d}\nu.
 \end{split}
\]
All integrals are well defined: each function on the left of a submean inequality is bounded above, and its negative part is controlled by that of its integrable unsmoothed counterpart plus a constant.

Radial averages of a plurisubharmonic function decrease when the radius decreases. Therefore $S_{a,j}(H_j+q_a)$ decreases to $z+q_a$ on the support of $\eta_a$. To identify the limit, fix an index $k$, bound $H_j\leq H_k$ for $j\geq k$, let $j\to\infty$ in the radius, and then let $k\to\infty$. The reverse inequality is the submean inequality. The nonnegative difference in the last display is dominated by an integrable constant minus $z$ on this compact support. Thus its integral tends to zero. The final integral tends to zero by $L^1(\nu)$ convergence. The subsequence criterion proves \eqref{eq:convolution-moving-integrals} for the original sequence.

The exponential is Lipschitz on the common upper range, so its moving $L^1$ error also tends to zero. If $\zeta>0$ is smooth, then $\zeta \mathrm{e}^{\lambda z}=\mathrm{e}^{\lambda z+\log\zeta}$ is bounded and quasi-plurisubharmonic. For any fixed bounded quasi-plurisubharmonic function $F$, the same local submean argument and dominated convergence in \eqref{eq:measure-smoothing-duality} give
\[
 \int_XF\,\mathrm{d}\nu_j\longrightarrow\int_XF\,\mathrm{d}\nu.
\]
Apply this to $F=\zeta \mathrm{e}^{\lambda z}$. Subtracting two positive smooth tests proves the required convergence for every smooth test function. Uniformly bounded masses and uniform approximation of continuous functions by smooth ones prove weak convergence.
\end{proof}

\Needspace{25\baselineskip}


\section{Curvature of line bundles}

Let $M$ be a complex manifold and let $\theta$ be a smooth real $(1,1)$-form on $M$. To relate potentials to holomorphic sections, we ask when a current $\theta+\ddc u$, with $u$ locally integrable, can be the curvature current of a metric on a line bundle. The transition law for the local weights of such a metric forces a closedness condition on the background.

\begin{proposition}\label{prop:chern-closed}\index{Line bundle!curvature}
The curvature current of a singular Hermitian metric with locally integrable weights on a holomorphic line bundle is closed. In particular, a nonclosed form $\theta$ cannot itself be the curvature current of such a metric, and $\theta+\ddc u$ is not closed on any neighborhood of a point where $\mathrm{d}\theta\neq0$.
\end{proposition}
\begin{proof}
Let $h$ be a singular Hermitian metric on a holomorphic line bundle $L\to M$. Choose a local holomorphic frame $e$ of $L$ and write its squared norm as $|e|_h^2=\mathrm{e}^{-\varphi}$; the curvature current is then locally $\ddc\varphi$. If $e'=ge$ is another frame, with $g$ holomorphic and nowhere zero, the new weight is $\varphi-\log|g|^2$. Since $\log|g|^2$ is pluriharmonic, these local currents agree on overlaps and glue to a global current, which is closed because $\mathrm{d}\ddc=0$. The final assertion follows from the identity
\[
 \mathrm{d}(\theta+\ddc u)=\mathrm{d}\theta
\]
of currents.
\end{proof}

Accordingly, when we study sections subject to multiplier-ideal conditions, we specify both the line bundle, whose curvature is closed, and any auxiliary Hermitian form on $M$. An algebraic interpretation then requires relating the invariants of the line bundle to this analytic background.

\section{Polar curvature}
\label{sec:psh-polar-invariance}

Bounded changes preserve more than logarithmic coefficients: they preserve the whole curvature measure on the common pole set. A local energy estimate puts the bounded difference in a Sobolev space, while pole sets admit smooth tests of arbitrarily small energy. These tests exclude every coefficient measure of its Hessian from the pole set. Polar curvature is closed even when the pole set is not closed. A local elliptic equation therefore recovers it from its scalar Hermitian trace.

\begin{lemma}\label{lem:bounded-measure-laplacian-energy}
Let $U\subset\mathbb R^d$ be open and let $w\in L^\infty_{\mathrm{loc}}(U)$. If $\Delta w$ is a locally finite signed Radon measure, then $w\in W^{1,2}_{\mathrm{loc}}(U)$.
\end{lemma}
\begin{proof}
Fix $0\leq\psi\leq1$ smooth and compactly supported on $U$, and a relatively compact open neighborhood $V$ of its support. Mollify $w$ inside $V$, writing $w_\varepsilon=w*\rho_\varepsilon$ and $\nu_\varepsilon=(\Delta w)*\rho_\varepsilon$. For small $\varepsilon$, integration by parts and Young's inequality give
\[
 \begin{split}
 \int\psi^2|\nabla w_\varepsilon|^2
 &=-\int\psi^2w_\varepsilon\nu_\varepsilon
       -2\int\psi w_\varepsilon\nabla\psi\cdot\nabla w_\varepsilon\\
 &\leq \|w\|_{L^\infty(V)}|\Delta w|(V)
       +\tfrac12\int\psi^2|\nabla w_\varepsilon|^2
       +2\int w_\varepsilon^2|\nabla\psi|^2.
 \end{split}
\]
The right side bounds the local gradient norm independently of $\varepsilon$. Taking $\psi=1$ on any prescribed smaller compact set, weak Sobolev compactness and convergence of mollifications in $L^2$ identify the weak limit with $w$. This proves the assertion.
\end{proof}

\begin{lemma}\label{lem:pole-small-energy-tests}\index{Pluripolar set!smooth energy tests}
Let $v$ be quasi-plurisubharmonic on a complex manifold, let $U$ be a coordinate ball with compact closure in a larger chart, and let $K\subset U\cap\{v=-\infty\}$ be compact. For every open $O\subset U$ containing $K$, there are real $\varphi_j\in C^\infty_c(O)$ with $0\leq\varphi_j\leq1$, equal to one near $K$, such that $\|\varphi_j\|_{W^{1,2}(U)}\to0$. The norm uses the Euclidean coordinates.
\end{lemma}
\begin{proof}
A smooth quadratic correction and subtraction of a constant give a negative plurisubharmonic function $h$ on $U$ with the same pole set as $v$. Put $h_M=\max(h,-M)$ for $M\geq1$ and choose $0\leq\psi\leq1$ in $C^\infty_c(O)$, equal to one near $K$. On a fixed compact neighborhood of $\operatorname{supp}\psi$, the positive measures $\Delta h_M$ have uniformly bounded mass: choose a nonnegative smooth compactly supported $\chi$ on $U$ equal to one there and use
\[
 \int\chi\,\Delta h_M=\int h_M\Delta\chi,
 \qquad |h_M|\leq|h|\in L^1_{\mathrm{loc}}(U).
\]
Apply the preceding integration-by-parts estimate to mollifications of $h_M$. Their absolute values are at most $M$, and their squares are at most $M$ times their absolute values. The uniform Laplacian mass and local $L^1$ bounds therefore give, after passing to the weak Sobolev limit,
\[
 \int\psi^2|\nabla h_M|^2\leq CM,
 \qquad
 \left\|\frac{-\psi h_M}{M}\right\|_{W^{1,2}(U)}^2\leq\frac{C}{M}.
\]
Here $C$ depends on the fixed cutoffs and the local $L^1$ norm of $h$, but not on $M$. The second inequality also uses $h_M^2\leq M|h|$ in the terms containing $\psi$ and $\nabla\psi$.

The Sobolev function $-\psi h_M/M$ lies between zero and one. It equals one on a neighborhood of $K$, since $h$ is upper semicontinuous and $\{h<-M\}$ is an open neighborhood of $K$. Extend it by zero and mollify with sufficiently small radius. Its support stays in $O$, it remains equal to one on a smaller neighborhood of $K$, and its Sobolev norm does not increase. Taking $M\to\infty$ gives the required smooth functions.
\end{proof}

Restrictions below are taken coefficientwise on the Radon measures of a current of order zero. They therefore commute with smooth changes of coordinates.

\Needspace{10\baselineskip}
\begin{proposition}\label{prop:bounded-type-polar-invariance}\index{Singularity type!polar curvature invariance}\index{Polar curvature!bounded changes}
Let $u,v$ be quasi-plurisubharmonic on a complex manifold and locally equivalent up to bounded constants. Their pole sets agree, say $S=\{u=-\infty\}=\{v=-\infty\}$, and
\[
 \mathbf1_S\ddc u=\mathbf1_S\ddc v.
\]
Consequently, for any smooth real $(1,1)$-forms $\theta,\beta$,
\[
 \mathbf1_S(\theta+\ddc u)=\mathbf1_S(\beta+\ddc v).
\]
\end{proposition}
\begin{proof}
The two local order inequalities give the same pole set. Off $S$ put $w=u-v$ and assign arbitrary bounded values on $S$. This defines the locally bounded distribution $u-v$, because $S$ has Lebesgue measure zero. The coefficient measures of $\ddc w$ are locally finite: after a smooth correction, each Hessian of $u$ and $v$ is a positive current of order zero. Its Euclidean trace makes $\Delta w$ a signed measure. Thus \cref{lem:bounded-measure-laplacian-energy} gives $w\in W^{1,2}_{\mathrm{loc}}$.

Each real and imaginary part of a coefficient measure of $\ddc w$ is consequently a continuous functional on $W^{1,2}_0$ on a smaller ball, since a second distributional derivative of $w$ is a first derivative of its $L^2$ gradient. Let $\nu$ be any one of these real signed measures, with Jordan parts $\nu^+,\nu^-$. Choose a compact $K$ in the intersection of $S$ with a positive Hahn set, so $\nu^-(K)=0$. For $\epsilon>0$, take an open neighborhood $O$ of $K$ in the ball with $\nu^-(O)<\epsilon$. The tests of \cref{lem:pole-small-energy-tests}, applied to $u$, satisfy
\[
 \nu^+(K)\leq\int\varphi_j\,\mathrm d\nu^+
       =\langle\nu,\varphi_j\rangle
          +\int\varphi_j\,\mathrm d\nu^-
       \leq|\langle\nu,\varphi_j\rangle|+\epsilon.
\]
Continuity on the Sobolev space makes the first term tend to zero. Letting $\epsilon\downarrow0$ and using inner regularity gives $\nu^+(S)=0$ locally. Applying the same argument to $-\nu$ gives $\nu^-(S)=0$. All coefficient measures of $\ddc w$ therefore vanish on $S$, proving the first identity. Smooth forms also assign zero mass to $S$, which proves the second.
\end{proof}

For a closed positive current, local potentials turn equality of scalar traces into an elliptic equation. If both currents are carried by a Lebesgue null set, the resulting smooth difference vanishes.

\Needspace{10\baselineskip}
\begin{lemma}\label{lem:closed-null-current-trace}\index{Closed positive current!trace on a null set}
Let $T,S$ be closed positive $(1,1)$-currents on a complex manifold of dimension $n\geq1$, both carried coefficientwise by a Borel set of Lebesgue measure zero. For any Hermitian form $\eta$,
\[
 T\wedge\eta^{n-1}=S\wedge\eta^{n-1}
 \quad\Longrightarrow\quad T=S.
\]
\end{lemma}
\begin{proof}
On a smaller coordinate ball, take plurisubharmonic local potentials $a,b$ for $T,S$, by \cite[Chapter III, Proposition 1.19, p.~135]{DemCADG}, with the normalization of this book. Put $h=a-b\in L^1_{\mathrm{loc}}$. The trace identity gives $\ddc h\wedge\eta^{n-1}=0$ in distributions. Dividing by the smooth positive volume form $\eta^n$ gives a scalar second-order operator whose principal matrix is a positive multiple of the inverse Hermitian matrix of $\eta$. It is uniformly elliptic on every smaller ball, with smooth coefficients. Local elliptic regularity therefore makes $h$ smooth. Thus $T-S=\ddc h$ is smooth and carried by a Lebesgue null set. Its coefficient measures must vanish, proving equality locally and hence globally.
\end{proof}

Proper pushforward commutes with coefficientwise restriction to the inverse image of a Borel set. This gives exact compatibility of polar curvature under modifications, even when its pole set is not closed.

\Needspace{12\baselineskip}
\begin{proposition}\label{prop:polar-curvature-pushforward}\index{Polar curvature!pushforward}
Let $\pi:Y\to X$ be a modification of connected compact complex manifolds, let $\theta$ be a smooth real $(1,1)$-form on $X$, and let $u\in\PSH(X,\theta)$. Put $S=\{u=-\infty\}$ and
\[
 T=\theta+\ddc u,\qquad
 T_Y=\pi^*\theta+\ddc(u\circ\pi).
\]
Then
\begin{equation}\label{eq:polar-curvature-pushforward}
 \pi_*\bigl(\mathbf1_{\pi^{-1}(S)}T_Y\bigr)=\mathbf1_ST.
\end{equation}
\end{proposition}
\begin{proof}
Pullback is genuine and admissible by \cref{prop:psh-descent-modification}. Change of variables on the isomorphism locus, whose complement has Lebesgue measure zero on both spaces, gives $\pi_*(u\circ\pi)=u$ as distributions and $\pi_*\pi^*\theta=\theta$ as currents. Proper pushforward commutes with $\ddc$, so $\pi_*T_Y=T$.

For any Borel $B\subset X$, the coefficient measures of $\pi_*T_Y$ are pushforwards of the coefficient measures of $T_Y$ multiplied by the smooth coefficients arising from pullback of a test form. Restriction therefore gives $\pi_*(\mathbf1_{\pi^{-1}(B)}T_Y)=\mathbf1_B\pi_*T_Y$. Apply this to the Borel pole set $S$. No closedness of its restricted current is used.
\end{proof}

The closedness of polar curvature follows from the bounded cutoff argument for complete pluripolar restrictions. The cutoff approaches the complement of the pole set pointwise, so the pole set need not be closed.

\Needspace{12\baselineskip}
\begin{proposition}\label{prop:polar-curvature-closedness}\index{Polar curvature!closedness}
Let $X$ be a complex manifold of dimension $n\geq1$, let $u$ be quasi-plurisubharmonic on $X$, and let $\theta$ be a smooth real $(1,1)$-form. Then
\[
 \operatorname{Pol}_\theta(u)
       :=\mathbf1_{\{u=-\infty\}}(\theta+\ddc u)
\]
is a closed positive $(1,1)$-current.
\end{proposition}
\begin{proof}
Locally, a smooth quadratic correction and subtraction of a constant give a negative plurisubharmonic function $a$ with pole set $S=\{u=-\infty\}$. Put $\Theta=\ddc a$. Since smooth forms have no coefficient mass on $S$, the polar current equals $\mathbf1_S\Theta$ and is positive. In dimension one it is closed by degree.

Suppose $n\geq2$. Choose a smooth convex nondecreasing function $\chi$ on $\mathbb R$ with $\chi(t)=0$ for $t\leq1/2$ and $\chi(1)=1$, and a smooth function $0\leq\vartheta\leq1$ which is zero near zero and one near one. Set
\[
 w_k=\chi(e^{a/k}),\qquad k\geq1.
\]
These functions are bounded plurisubharmonic, with $0\leq w_k\leq1$. They vanish on $S$ and tend to one at every point outside $S$. Hence $\vartheta(w_k)\Theta$ converges coefficientwise locally in total variation to $\mathbf1_{X\setminus S}\Theta$. Also $\vartheta'(w_k)$ vanishes on $S$ and tends to zero everywhere.

Fix a constant strongly positive $(n-2,n-2)$-form $\beta$ and put $\Xi=\Theta\wedge\beta$. The bounded Bedford--Taylor chain rule gives
\[
 \mathrm d\bigl(\vartheta(w_k)\Xi\bigr)
       =\vartheta'(w_k)\,\mathrm dw_k\wedge\Xi.
\]
Indeed, decrease smooth plurisubharmonic functions to $w_k$ on a smaller ball. Dominated convergence gives convergence of their $\vartheta$-weighted currents. The bounded Bedford--Taylor energy convergence and Cauchy--Schwarz give convergence of the differentiated expressions, proving the chain rule for $w_k$.
For a smooth compactly supported test one-form $\alpha$, choose a nonnegative smooth cutoff $\tau$ equal to one near its support. Since $w_k\geq0$ and $\Xi$ is closed and positive, bounded integration by parts gives
\[
 2\int\tau\,\mathrm dw_k\wedge\dc w_k\wedge\Xi
       \leq\int\tau\,\ddc(w_k^2)\wedge\Xi
       =\int w_k^2\ddc\tau\wedge\Xi\leq C.
\]
Here $C$ depends only on the local mass of $\Xi$ and the fixed cutoff. The Cauchy--Schwarz inequality for a positive current now yields
\[
 \bigl|\langle\vartheta'(w_k)\,\mathrm dw_k\wedge\Xi,
                    \alpha\rangle\bigr|^2
       \leq C_\alpha\int_{\operatorname{supp}\alpha}
                 |\vartheta'(w_k)|^2\,\Xi\wedge\kappa
       \longrightarrow0,
\]
where $\kappa$ is the Euclidean Hermitian form. The limit follows by dominated convergence against the fixed locally finite trace measure. Thus $\mathbf1_{X\setminus S}\Theta\wedge\beta$ is closed. Constant strongly positive forms span the needed test forms, so $\mathbf1_{X\setminus S}\Theta$ is closed. Subtracting it from $\Theta$ proves closedness of $\mathbf1_S\Theta$. The local currents agree on overlaps and give the asserted global current.
\end{proof}

